% =============================================================================
%  main.tex - preambulo, portada y \input de las secciones del reporte final.
%
%  ESTE ARCHIVO ES DE TODO EL EQUIPO. Solo se toca para agregar una seccion
%  (un \input y su \addbibresource). El texto va en secciones/<seccion>.tex.
%
%  Compilar desde esta carpeta:   latexmk -pdf main.tex
% =============================================================================

% cas-dc = plantilla CAS de Elsevier a dos columnas. A diferencia de
% elsarticle no carga natbib, asi que biblatex entra sin opciones extra.
% Ya carga graphicx, amsmath, booktabs, multirow, xcolor, hyperref y la
% tipografia STIX; no hay que volver a pedirlos.
% longmktitle (portada de mas de una pagina) truena en TeX Live 2025: usa una
% funcion de expl3 que ya se retiro. La portada cabe en una pagina sin ella.
\documentclass[a4paper,fleqn]{cas-dc}

% --- Idioma ------------------------------------------------------------------
% Los ajustes de babel-spanish van en \spanishoptions ANTES de cargar babel:
%   es-tabla        -> "Tabla 1" en vez de "Cuadro 1"
%   es-nodecimaldot -> punto decimal en matematicas (12.46 y no 12,46)
%   es-noquoting    -> evita conflictos de comillas con otros paquetes
\def\spanishoptions{es-tabla,es-nodecimaldot,es-noquoting}
\usepackage[spanish]{babel}
\usepackage{csquotes}

% cas-dc escribe "ABSTRACT", "ARTICLE INFO" y "Keywords" fijos en ingles.
% babel cambia \abstractname al entrar al documento, asi que se reescribe
% en \captionsspanish; el recuadro de palabras clave se redefine completo.
\addto\captionsspanish{\def\abstractname{R\,E\,S\,U\,M\,E\,N}}
\makeatletter
\ExplSyntaxOn
\RenewDocumentEnvironment { keywords } { O{ Palabras~clave } }
   {
     \tex_global:D \tex_setbox:D \g_stm_key_box = \vtop \bgroup
     \hsize=.25 \textwidth
     \cs_set:Nn \sep: { \par }
     \cs_set_eq:NN \sep \sep:
     \parindent \z@
     I\,N\,F\,O\,R\,M\,A\,C\,I\,Ó\,N \par \skip_vertical:n { -3pt }
     \rule{.25 \textwidth}{.2pt}\par\footnotesize
     \noindent \textit { #1 }:  \par
   }
   { \egroup }

% Pie de pagina: sin "Preprint submitted to Elsevier" (es un reporte de
% clase) y con "Página n de m".
\cs_set:Npn \__first_footerline:
  { \group_begin: \small \sffamily \__short_authors: \group_end: }
\cs_set:Npn \__first_foot:
  {
    \parbox[t]{\textwidth}
      { \rule{\textwidth}{.2pt}\\ \__first_footerline: \hfill
        Página~\thepage{}~de~\lastpage }
  }
\cs_set:Npn \__cas_foot:
  {
    \parbox[t]{\textwidth}
      { \rule{\textwidth}{.2pt}\\ \sffamily\small \__first_footerline:
        \hfill Página~\thepage{}~de~\lastpage }
  }
% \ps@cas copia el pie al cargarse la clase; se reactiva para tomar el nuevo.
\ps@cas
% Sin la nota vacia "ORCID(s):" al pie de la primera pagina.
\cs_set:Npn \printorcid { }
% "y" en lugar de "and" antes del ultimo autor (copia de \stmauthors de
% cas-common.sty con ese unico cambio).
\RenewDocumentCommand \stmauthors { }
   {
     \group_begin:
     \stmAuthorSetup { type = authors }
     \l_stm_au_setup_tl
     \par \vskip\l_stm_augroup_before_dim
     \l_stm_augroup_align_tl
     \l_stm_augroup_size_tl
     \l_stm_augroup_shape_tl
     \l_stm_augroup_weight_tl
     \color{ \l_stm_augroup_color_tl }
     \bool_if:NTF \g_stm_augr_bool
       { \seq_use:Nn \g_stm_augr_seq { \par } }
       {
         \seq_use:cnnn { g_stm_au\int_use:N\g_stm_augr_int _seq }
         { ~y~ } { ,~ } { ~y~ }
       }
     \par\vskip\l_stm_augroup_after_dim
     \group_end:
   }
\ExplSyntaxOff
\makeatother

% --- Tipografia y espaciado --------------------------------------------------
\usepackage{microtype}
% A dos columnas una palabra larga sin guiones se sale del margen; esto deja
% que LaTeX estire un poco los espacios antes de desbordar.
\emergencystretch=2em

% --- Matematicas -------------------------------------------------------------
\usepackage{amsthm}
\theoremstyle{definition}
\newtheorem{definition}{Definición}

% --- Tablas, figuras y listas ------------------------------------------------
\usepackage{tabularx}
\usepackage{ragged2e}
\usepackage{enumitem}
\usepackage{tikz}             % figuras/clave*.tex (scripts/figura_clave.py)
\setlist{noitemsep, topsep=3pt, partopsep=0pt}
\graphicspath{{figuras/}{../figures/}}

% --- Bibliografia (biblatex + biber) -----------------------------------------
% Cada seccion trae su propio .bib (no se comparten) y todas se imprimen en
% una sola lista al final, como en los reportes de las etapas 1 y 2.
\usepackage[backend=biber,style=apa,sorting=nyt,maxbibnames=99]{biblatex}
\renewcommand*{\bibfont}{\small}
\setlength{\bibitemsep}{0pt}
\setlength{\bibhang}{1em}

% Etapa 1
\addbibresource{secciones/introduccion.bib}
\addbibresource{secciones/trabajo-relacionado.bib}
\addbibresource{secciones/marco-estocastico.bib}
% Etapa 2
\addbibresource{secciones/modelo-compra.bib}
\addbibresource{secciones/modelo-espacial.bib}
% Etapa 3

% --- Referencias cruzadas (despues de hyperref, que ya cargo la clase) -------
\hypersetup{linkcolor=blue,citecolor=blue,urlcolor=blue}
\usepackage[spanish]{cleveref}
% Secciones en romanos, como en la rubrica (I a V); las subsecciones quedan
% como IV.1, IV.1.1, etc.
\renewcommand{\thesection}{\Roman{section}}
\crefname{section}{sección}{secciones}
\Crefname{section}{Sección}{Secciones}
\crefname{table}{tabla}{tablas}
\Crefname{table}{Tabla}{Tablas}
\crefname{equation}{ecuación}{ecuaciones}
\Crefname{equation}{Ecuación}{Ecuaciones}

\begin{document}
\let\WriteBookmarks\relax
\def\floatpagepagefraction{1}
\def\textpagefraction{.001}

\shorttitle{Modelación estocástica para la reforestación}
\shortauthors{Luna, Alvarez, Sánchez y de la Peña}

\title[mode=title]{Compra bajo incertidumbre y distribución espacial de especies
  para la reforestación de 30 polígonos en San Luis Potosí}

\author[tec]{Angel Eduardo Luna}
\author[tec]{Irasema Alvarez}
\author[tec]{Carolina Sánchez Meneses}
\author[tec]{Andrea Estefanía de la Peña}

% Nota al pie sin marca en la portada con el enlace al codigo.
\nonumnote{Código y datos: \url{https://github.com/Lunaaaalj/Stochastic_RefInit}}

\affiliation[tec]{organization={Instituto Tecnológico y de Estudios Superiores
  de Monterrey, MA2004B, grupo 202, equipo 5},
  country={México}}

\begin{abstract}
La Comisión Nacional Forestal planea reforestar 30 polígonos (183.01~ha) del
matorral desértico de San Luis Potosí con diez especies nativas, al tresbolillo
y con una meta de 658 plantas por hectárea. Se propone un modelo de dos
niveles. El primero decide cuántas plantas de cada especie comprar en cada
polígono para alcanzar la meta con probabilidad de al menos $1-\alpha$, al
menor costo y dentro del presupuesto: la restricción de probabilidad se reduce
al cuantil de una binomial negativa, estimado por simulación Monte Carlo, y la
elección de metas es una mochila de elección múltiple. El segundo acomoda las
plantas en un mosaico periódico de 52 posiciones para minimizar la competencia
entre vecinas, con un algoritmo memético validado contra un modelo de
programación entera. Con $\alpha=0.05$, el presupuesto de \$12{,}500/ha no
alcanza en ningún polígono: se requieren en promedio \$14{,}496/ha, y aceptar
más riesgo apenas reduce esa cifra, mientras que el margen de ajuste por
especie es la principal palanca de costo. Los 30 polígonos se reducen a seis
claves de plantación que disminuyen la competencia entre 14.5\,\% y 16.3\,\%
respecto al azar, a no más de 4.5\,\% del óptimo, y que superan a la secuencia
vigente, la cual, repetida fila tras fila, compite más que una plantación al
azar. En 1{,}000 siembras simuladas de una hectárea, la clave reduce en 18\,\%
la competencia esperada al año, y cumplir cada especie con 95\,\% de
probabilidad deja en 68\,\% la probabilidad de cumplirlas todas a la vez.
\end{abstract}

\begin{keywords}
reforestación \sep tresbolillo \sep restricciones de probabilidad \sep
simulación Monte Carlo \sep coloración de grafos \sep algoritmo memético
\end{keywords}

\maketitle

% --- Secciones ---------------------------------------------------------------
% Una seccion = un archivo. El orden sigue la rubrica: introduccion y
% justificacion, problematica y objetivo, modelacion y validacion,
% resultados, discusion y conclusiones.

% Orden de la rubrica: I Introduccion, II Marco teorico, III Descripcion del
% problema, IV Metodologia de solucion, V Experimentacion y resultados, y
% conclusiones. marco-teorico, metodologia y experimentacion solo abren su
% seccion e incluyen las subsecciones de los demas archivos.
\section{Introducción}
\label{sec:introduccion}

\subsection{Reforestación en zonas áridas}
\label{subsec:reforestacion}

La deforestación por cambio de uso de suelo es una de las principales causas de pérdida de biodiversidad y de servicios ecosistémicos \autocite{geist2002}. La restauración ecológica responde a este problema asistiendo la recuperación de un ecosistema degradado, dañado o destruido \autocite{ser2004}, y la reforestación es una de sus herramientas: establecer vegetación, mediante siembra o plantación, en terrenos que la tenían y la perdieron \autocite{conafor2010}. Su objetivo no es solo recuperar la cobertura vegetal, sino también la retención de suelo, la infiltración de agua, el hábitat para la fauna y la captura de carbono \autocite{chazdon2008}. Cuando la regeneración natural no basta, se recurre a la restauración activa, que planta especies seleccionadas \autocite{holl2011}.

Una reforestación se planea en tres etapas: el diagnóstico del sitio y la selección de especies, densidad y diseño de plantación; la ejecución, que incluye la obtención de la planta en vivero y su plantación en la temporada de lluvias; y el mantenimiento, que abarca la reposición de plantas muertas y la medición de la supervivencia \autocite{conafor2010}. Se da prioridad a especies nativas adaptadas al clima y al suelo local, en proporción a su frecuencia en la vegetación de la región. De los diseños de plantación, el tresbolillo desfasa las filas para formar triángulos equiláteros, lo que aprovecha mejor el espacio que el marco real y protege mejor contra el viento y la erosión.

\subsection{Competencia y facilitación entre plantas}
\label{subsec:competencia}

La competencia ocurre cuando plantas vecinas usan los mismos recursos limitados (agua, luz y nutrientes del suelo) y el consumo de una reduce lo disponible para la otra \autocite{craine2013}. Es más intensa entre plantas cercanas, de tamaño parecido y que toman los recursos de la misma forma, a la misma profundidad de raíz y en la misma época. En zonas áridas el agua es el recurso más limitante, por lo que plantar juntas especies con alta demanda de agua aumenta la competencia \autocite{padilla2006}.

Las interacciones no son solo negativas. Las plantas nodriza crean bajo su copa un microambiente con más sombra, humedad y nutrientes que favorece el establecimiento de otras especies \autocite{flores2003}. El mezquite (\textit{Prosopis}) es una nodriza común en el norte de México, y colocar especies sensibles bajo nodrizas es una estrategia reconocida para restaurar ambientes degradados \autocite{padilla2006}. En consecuencia, el acomodo de las especies, y no solo su número, influye en la supervivencia de la plantación.

Una reforestación mal planeada puede acomodar las plantas de forma que compitan entre sí, lo que reduce la supervivencia y deja sin cumplir la meta \autocite{conafor2010,holl2011}; reducirla a pocas especies tampoco es una salida, porque una plantación casi monoespecífica es más vulnerable a plagas, enfermedades y eventos climáticos extremos \autocite{liu2018}. A esto se suma un problema logístico: si se compran plantas de más se gasta dinero de forma innecesaria, y si se compran de menos no se alcanza la meta.

\subsection{El reto}
\label{subsec:reto}

El reto se ubica en la zona de la Línea de Transmisión Eléctrica Dominica--Charcas, en San Luis Potosí, dentro de la subprovincia Llanos y Sierras Potosino--Zacatecanos. La vegetación es principalmente matorral desértico rosetófilo y micrófilo, con un deterioro notable por sobrepastoreo. El socio formador, la Comisión Nacional Forestal (CONAFOR), planea reforestar 30 polígonos de entre 1.28 y 8~ha (183.01~ha en total) con diez especies nativas de agaves, nopales, mezquite y yuca, en un diseño al tresbolillo elegido por los fuertes vientos de la zona y con una meta de 658 plantas por hectárea.

El problema tiene dos partes. La primera es estimar cuántas plantas de cada especie comprar, sabiendo que en cada polígono ya existen plantas de regeneración natural y que no todas las plantas nuevas sobreviven. La segunda es decidir cómo distribuir las especies dentro de cada polígono para reducir la competencia entre plantas vecinas, respetando la cantidad comprada de cada especie.

\subsection{Justificación y objetivo}
\label{subsec:objetivo}

Una buena distribución reduce la competencia y aprovecha la facilitación, lo que aumenta la supervivencia; en una zona árida, cada planta perdida es difícil de reponer. Por su parte, estimar bien la cantidad a comprar evita tanto el gasto innecesario como el incumplimiento de la meta. Un método sistemático para ambas decisiones permite a CONAFOR planear con datos en lugar de depender de un orden de plantación fijo.

El objetivo general es proponer un modelo matemático y computacional que permita a CONAFOR:
\begin{enumerate}
  \item determinar, para cada polígono, cuántas plantas de cada especie adquirir para alcanzar la meta de plantación con una probabilidad prefijada, al menor costo y dentro del presupuesto, considerando la incertidumbre en la supervivencia de las plantas nuevas; y
  \item distribuir esas plantas en la retícula al tresbolillo de modo que se minimice la competencia entre especies vecinas, con un arreglo que pueda llevarse a campo.
\end{enumerate}

\section{Marco teórico}
\label{sec:marco-teorico}

El marco teórico tiene dos partes. La primera revisa los trabajos que asignan elementos a posiciones para reducir la competencia entre vecinos y el marco de coloración de grafos que los unifica, base del modelo espacial. La segunda reúne las herramientas de probabilidad con que se modela la compra bajo incertidumbre: las distribuciones de conteo, el método de Monte Carlo y las cadenas de Markov.

\section{Trabajo relacionado}
\label{sec:trabajo-relacionado}

La distribución de especies pertenece a una familia de problemas que asignan elementos a posiciones de modo que se reduzca la competencia, la interferencia o la incompatibilidad entre vecinos. Esta sección revisa cinco trabajos de esa familia (\cref{subsec:competencia-literatura}) y el marco de coloración de grafos que los unifica (\cref{subsec:coloracion}).

\subsection{Reducción de competencia}
\label{subsec:competencia-literatura}

El trabajo más cercano al reto es el de \textcite{elizalde2025}, que asignan diez especies a las posiciones de siembra de un polígono siguiendo los lineamientos de CONAFOR (625 plantas por hectárea). El terreno es un grafo cuyos vértices son posiciones y cuyos arcos indican vecindad; una variable binaria indica qué especie ocupa cada posición, y un coeficiente de competencia por par de especies se contabiliza solo cuando ambas quedan en posiciones vecinas. Resuelven un modelo entero con CBC y un algoritmo genético de cromosoma entero. Con un límite de tres horas, el modelo exacto obtuvo soluciones factibles sin garantía de calidad entre 169 y 324 nodos y ninguna con 361, 400 y 625 nodos, mientras que el genético resolvió todos los tamaños (349.44~s en el mayor) con calidad igual o muy cercana a la del exacto donde hubo comparación.

\textcite{eisenblatter2002} estudian la asignación de frecuencias en redes celulares, donde la interferencia entre transmisores cumple el papel de la competencia entre plantas: cada transmisor es un vértice, los canales son colores y una penalización por par de transmisores mide la interferencia. Formulan el problema de mínima interferencia como un programa entero, muestran que es fuertemente NP-difícil y reportan que las relajaciones lineales dan cotas débiles, mientras que heurísticas como DSATUR seguidas de búsqueda local mejoran mucho los planes reales. \textcite{akplogan2011} asignan cultivos a unidades de terreno a lo largo de varios años como un problema de satisfacción de restricciones ponderadas; con bloques interdependientes, la instancia de 120 unidades no se cerró en 48 horas. \textcite{liu2026} planean cultivos con una matriz de adyacencia $W$ entre unidades vecinas y una matriz de interacción $M$ entre cultivos, positiva para pares complementarios y negativa para pares competitivos, e incorporan incertidumbre en rendimientos y precios con escenarios Monte Carlo y optimización robusta. Por último, \textcite{kobal2025} muestran en bosques mixtos que los índices de competencia que consideran la distancia a los vecinos explican mejor el crecimiento que los que la ignoran (el $R^{2}$ ajustado pasa de 0.325 sin competencia a 0.647 con el índice de Hegyi).

\begin{table*}[t]
  \centering
  \caption{Comparación de los trabajos sobre reducción de competencia.}
  \label{tab:comparacion-competencia}
  \small
  \begin{tabularx}{\textwidth}{@{}>{\RaggedRight\arraybackslash}p{2.7cm}>{\RaggedRight\arraybackslash}p{3.1cm}*{2}{>{\RaggedRight\arraybackslash}X}@{}}
    \toprule
    Referencia & Dominio & Criterio de competencia o incompatibilidad & Método \\
    \midrule
    \parencite{elizalde2025} & Reforestación & Coeficiente por par de especies en posiciones vecinas & Programación entera y algoritmo genético \\
    \parencite{eisenblatter2002} & Telecomunicaciones & Interferencia por par de transmisores (0 a 1) & Programación entera, heurísticas y cotas \\
    \parencite{akplogan2011} & Agricultura & Efecto del cultivo previo, compatibilidad con el suelo y agrupación espacial & Restricciones ponderadas con ramificación y acotamiento \\
    \parencite{liu2026} & Agricultura bajo incertidumbre & Matriz de interacción entre cultivos vecinos & Optimización robusta con escenarios Monte Carlo \\
    \parencite{kobal2025} & Manejo forestal & Índices según tamaño y distancia de los vecinos & Regresión con optimización de umbrales \\
    \bottomrule
  \end{tabularx}
\end{table*}

La \cref{tab:comparacion-competencia} resume los cinco trabajos. De ellos se desprenden tres lecciones para el reto. Primero, la forma dominante de medir la competencia es un coeficiente por par de elementos vecinos en un grafo \parencite{elizalde2025,eisenblatter2002,liu2026}, y la distancia importa \parencite{kobal2025}, lo que respalda limitar la competencia a los vecinos inmediatos de la retícula. Segundo, los modelos exactos pierden escalabilidad rápidamente \parencite{elizalde2025,akplogan2011}, lo que justifica las metaheurísticas \parencite{elizalde2025,eisenblatter2002}; un modelo exacto sigue siendo útil en instancias pequeñas para medir qué tan lejos queda la heurística del óptimo. Tercero, solo \textcite{liu2026} incorporan incertidumbre, y ninguno estima la cantidad de plantas a adquirir cuando la supervivencia es incierta, que es la primera parte del reto.

\subsection{Coloración de grafos con penalización}
\label{subsec:coloracion}

Sea $G=(V,E)$ un grafo no dirigido. Una coloración propia con $k$ colores es una función $c:V\to\{1,\dots,k\}$ tal que $c(u)\neq c(v)$ para toda arista $uv\in E$; el menor $k$ para el que existe es el número cromático $\chi(G)$, y encontrarlo es NP-difícil \parencite{eisenblatter2002}. Como la coloración propia no basta para modelar requisitos prácticos, existen generalizaciones que, en lugar de prohibir los conflictos, minimizan una penalización total sobre las aristas \parencite{eisenblatter2002}, y sus variantes se han extendido a problemas de programación de tareas y horarios \parencite{dewerra2015}.

La distribución de especies se representa con un grafo cuyos vértices son las posiciones de siembra, cuyas aristas unen posiciones vecinas, cuyos colores son las especies y en el que cada arista pesa la competencia entre las especies asignadas a sus extremos. Esta representación coincide con la de \textcite{elizalde2025} y con la matriz de adyacencia de \textcite{liu2026}. En el tresbolillo, cada posición interior tiene seis vecinas equidistantes, de modo que el grafo es una retícula triangular y todas las aristas representan la misma distancia.

Una coloración propia no sirve directamente, por dos razones: no respeta la cantidad de plantas de cada especie, y trata la competencia como todo o nada, cuando en realidad es parcial y distinta para cada par de especies (incluido el par de una especie consigo misma). El modelo de mínima interferencia de \textcite{eisenblatter2002} y el de \textcite{elizalde2025} comparten, en cambio, la misma estructura: variables binarias de asignación, una restricción que asigna exactamente un color a cada vértice y una penalización por arista que se activa según los colores de sus extremos, cuya suma se minimiza. Esta coloración con penalización y conteo fijo por color es el marco que adopta la distribución espacial de este trabajo.

\subsection{Herramientas para modelar la incertidumbre}
\label{sec:marco-estocastico}

La compra se decide antes de saber cuántas plantas nuevas sobrevivirán. Esta sección reúne las distribuciones con que se modelan los conteos del problema y el método de Monte Carlo con el que se evalúan las decisiones bajo incertidumbre.

\subsubsection{Distribuciones de conteo}
\label{subsec:distribuciones}

\paragraph{Binomial} Si se plantan $x$ plantas de una especie y cada una sobrevive, de forma independiente, con probabilidad $q$, el número de sobrevivientes es $S(x)\sim\operatorname{Bin}(x,q)$, con
\begin{equation}
  \mathbb{P}(S(x)=k)=\binom{x}{k}q^{k}(1-q)^{x-k},\quad k=0,\dots,x
  \label{eq:binomial}
\end{equation}
y media $xq$ \autocite{:berger2002}. Comprar $\lceil r/q\rceil$ plantas para obtener $r$ sobrevivientes acierta solo en promedio: con esa compra, la meta se incumple con una probabilidad cercana a la mitad.

\paragraph{Binomial negativa} En una secuencia de ensayos de Bernoulli independientes con probabilidad de éxito $q$, el número de ensayos $X_r$ hasta obtener el $r$-ésimo éxito tiene distribución binomial negativa \autocite{:berger2002}:
\begin{equation}
  \mathbb{P}(X_r=x)=\binom{x-1}{r-1}q^{r}(1-q)^{x-r},\quad x=r,r+1,\dots
  \label{eq:binneg}
\end{equation}
En el problema, $X_r$ es el número de plantas que hay que plantar, una tras otra, hasta que $r$ sobrevivan. Las dos distribuciones se relacionan por la identidad
\begin{equation}
  \mathbb{P}\bigl(S(x)\geq r\bigr)=\mathbb{P}\bigl(X_r\leq x\bigr),
  \label{eq:dualidad}
\end{equation}
porque al menos $r$ de las primeras $x$ plantas sobreviven si y solo si la $r$-ésima sobreviviente aparece a más tardar en la planta $x$. En consecuencia, la menor compra que obtiene $r$ sobrevivientes con probabilidad al menos $1-\alpha$ es el cuantil $1-\alpha$ de $X_r$. Esta equivalencia convierte una restricción de probabilidad sobre la supervivencia en el cálculo de un cuantil.

\paragraph{Poisson} El número de ocurrencias de un evento en una región, cuando estas ocurren al azar con intensidad constante $\lambda$ por unidad de área, sigue una distribución de Poisson, $\mathbb{P}(E=k)=e^{-\lambda A}(\lambda A)^{k}/k!$ para una región de área $A$ \autocite{:berger2002}. Es el modelo natural para el número de plantas existentes por polígono si estas se dispersaron al azar. Su media y su varianza son iguales; cuando los conteos muestran una varianza mayor que la media (sobredispersión), la alternativa habitual es una binomial negativa vista como mezcla de Poisson.

\subsubsection{Método de Monte Carlo}
\label{subsec:montecarlo}

El método de Monte Carlo estima una cantidad $\theta=\mathbb{E}[g(\mathbf{X})]$ promediando evaluaciones del modelo sobre $K$ realizaciones independientes de la entrada aleatoria \autocite{:robert2004}:
\begin{equation}
  \hat{\theta}_K=\frac{1}{K}\sum_{k=1}^{K}g(\mathbf{X}_k)\xrightarrow{\text{c.s.}}\theta,
  \label{eq:montecarlo}
\end{equation}
donde la convergencia la garantiza la ley fuerte de los grandes números. Cada realización, o escenario, es un valor completo de las variables aleatorias del modelo; se obtiene transformando números pseudoaleatorios uniformes en la distribución deseada \autocite{:ross2013} y sobre él el modelo se evalúa de forma determinista. Una probabilidad es el caso $g=\mathbf{1}\{\cdot\}$, de modo que $\mathbb{P}(S(x)\geq r)$ se estima con la fracción de escenarios en que la compra $x$ alcanza la meta, y un cuantil se estima con el estadístico de orden correspondiente de los escenarios.

Si $\sigma^{2}=\operatorname{Var}[g(\mathbf{X})]<\infty$, el teorema central del límite da el intervalo aproximado de nivel $1-\alpha$
\begin{equation}
  \hat{\theta}_K\pm z_{\alpha/2}\frac{S_K}{\sqrt{K}},
  \label{eq:intervalo}
\end{equation}
con $S_K$ la desviación estándar muestral; para que el semiancho no exceda un error $\varepsilon$ se requieren $K\geq(z_{\alpha/2}S_K/\varepsilon)^{2}$ escenarios \autocite{:ross2013}. El error decrece como $K^{-1/2}$ sin importar la dimensión del problema: reducirlo a la mitad cuesta cuatro veces más escenarios. El método admite modelos no lineales arbitrarios y entrega la distribución completa de la salida, no solo su media, pero sus resultados no son mejores que las distribuciones supuestas para las entradas \autocite{:robert2004}.

\subsubsection{Cadenas de Markov}
\label{subsec:markov}

Un proceso $\{X_t\}_{t\geq0}$ con espacio de estados finito $\mathcal{X}$ es una cadena de Markov homogénea si la probabilidad de pasar al estado $j$ depende solo del estado actual \autocite{:ross2014}:
\begin{equation}
  \mathbb{P}(X_{t+1}=j\mid X_t=i,X_{t-1},\dots,X_0)=\mathbb{P}(X_{t+1}=j\mid X_t=i)=P_{ij}.
  \label{eq:markov}
\end{equation}
Las probabilidades de transición forman una matriz estocástica $\mathbf{P}$, con $P_{ij}\geq0$ y $\sum_jP_{ij}=1$, y si $\boldsymbol{\pi}_0$ es la distribución inicial, la distribución tras $t$ pasos es $\boldsymbol{\pi}_0\mathbf{P}^{t}$ \autocite{:norris1997}. Una distribución $\boldsymbol{\pi}$ con $\boldsymbol{\pi}\mathbf{P}=\boldsymbol{\pi}$ es estacionaria: describe la fracción de tiempo que la cadena pasa en cada estado.

En el problema hay dos cadenas naturales. La primera describe el orden de plantación: si cada planta es un estado, su especie, $P_{ij}$ es la probabilidad de que la siguiente planta del recorrido sea de la especie $j$ cuando la actual es de la $i$. Un arreglo que reduce la competencia tiene $P_{ii}=0$, porque nunca planta dos individuos de la misma especie seguidos, y su distribución estacionaria es la composición del arreglo. La segunda describe la vida de cada posición de la retícula, con estados vacía, plántula y establecida: una plántula muere o se establece, y una posición vacía solo cambia si se replanta. Con la distribución inicial dada por el conteo, esta cadena estimaría cuántas posiciones quedarán vacías en cada campaña de reposición. Este trabajo usa la primera para describir las claves (\cref{subsec:simulacion}) y se limita al horizonte de un año en la segunda, con la supervivencia binomial de la \cref{eq:binomial}.


\section{Descripción del problema}
\label{sec:problema}

\subsection{Decisiones, información y resultados esperados}
\label{subsec:decisiones}

El reto pide dos decisiones relacionadas pero distintas.
\begin{enumerate}
  \item \textbf{Cantidad a adquirir.} Cuántas plantas de cada una de las diez especies comprar en cada polígono, dado que ya existen plantas de regeneración natural y que una fracción de las plantas nuevas morirá.
  \item \textbf{Distribución de especies.} Qué especie ocupa cada posición de la retícula al tresbolillo, buscando reducir la competencia entre vecinos.
\end{enumerate}
La primera alimenta a la segunda: la compra fija cuántas plantas de cada especie hay que acomodar.

La información disponible es el plan de siembra (plantas por hectárea de cada especie, que suman 658); la superficie de cada uno de los 30 polígonos; los conteos de plantas existentes por especie en cada polígono; la supervivencia por especie medida de 2021 a 2023 en una muestra del 5\,\% de la superficie; el costo unitario de la planta de cada especie; y la secuencia de 22 posiciones (la clave de plantación) con la que actualmente se ordena la plantación en cada 100~m lineales. El plan admite que la cantidad de cada especie se desvíe de la proporción nominal dentro de un margen de entre $\pm$5\,\% y $\pm$10\,\%.

Los resultados esperados son una compra por especie y polígono que cumpla la meta del plan pese a la mortalidad, y un arreglo espacial de fácil implementación en campo que reduzca la competencia frente al arreglo actual.

\subsection{Variables, parámetros e información faltante}
\label{subsec:variables}

La \cref{tab:variables} resume los elementos que intervienen en el problema y cómo se tratan. Las variables de decisión son el núcleo de las dos preguntas del reto. La supervivencia de las plantas nuevas es la fuente de incertidumbre que justifica un enfoque estocástico en lugar de una regla determinista. La competencia entre especies no viene en los datos, pero es indispensable porque es justamente lo que la distribución espacial busca reducir. El agua, la disponibilidad en vivero y el terreno no son indispensables para una primera versión, pero condicionan qué tan realista es la solución y son los primeros candidatos para refinarla.

\begin{table*}[t]
\centering
\caption{Elementos del problema de reforestación y su tratamiento.}
\label{tab:variables}
\small
\begin{tabularx}{\textwidth}{@{}>{\RaggedRight\arraybackslash}p{0.25\textwidth}>{\RaggedRight\arraybackslash}p{0.16\textwidth}>{\RaggedRight\arraybackslash}X@{}}
\toprule
Elemento & Tipo & Descripción y tratamiento \\
\midrule
Cantidad a comprar por especie y polígono & Variable de decisión & Número de plantas nuevas a adquirir \\
Especie en cada posición de la retícula & Variable de decisión & Cuál de las diez especies ocupa cada sitio del tresbolillo \\
Plantas requeridas por especie & Parámetro conocido & 658 plantas/ha repartidas entre las diez especies según el plan, con un margen de ajuste por especie \\
Superficie de cada polígono & Parámetro conocido & De 1.28 a 8~ha; fija el número de plantas de la meta \\
Separación entre plantas & Parámetro conocido & Fijada por el tresbolillo y la densidad del plan, $\approx$4.19~m \\
Costo unitario por especie & Parámetro conocido & De \$18 a \$26.5 por planta, sin incluir la logística \\
Presupuesto & Parámetro conocido & \$12{,}500 por hectárea, fijado por el socio formador \\
Plantas existentes por especie y polígono & Conteo con incertidumbre & Se conoce un conteo por polígono, pero no la posición de cada planta \\
Supervivencia por especie & Variable aleatoria & Estimada con el monitoreo de 2021 a 2023; el número de sobrevivientes de una compra es incierto \\
Competencia entre pares de especies & Información a estimar & No viene en los datos del reto; se construye con literatura y criterio ecológico \\
Requerimientos de agua, disponibilidad en vivero y terreno por polígono & Información no disponible & Quedan fuera del modelo; se discuten como limitaciones \\
\bottomrule
\end{tabularx}
\end{table*}

\subsection{Requisitos del modelo}
\label{subsec:requisitos}

Hay dos criterios de optimización naturales: minimizar el costo de la compra, que responde a la primera decisión, y minimizar la competencia entre especies vecinas, sumada sobre todos los pares de posiciones adyacentes, que responde a la segunda. Ambos se conectan porque el margen de ajuste por especie permite cambiar la composición de cada polígono, y con ella tanto el costo como los pares de vecinos posibles. La \cref{tab:requisitos} enumera las condiciones que el modelo debe respetar.

\begin{table*}[t]
\centering
\caption{Condiciones que debe respetar el modelo.}
\label{tab:requisitos}
\small
\begin{tabularx}{\textwidth}{@{}>{\RaggedRight\arraybackslash}p{0.22\textwidth}>{\RaggedRight\arraybackslash}X>{\RaggedRight\arraybackslash}p{0.3\textwidth}@{}}
\toprule
Condición & Qué representa & Por qué se incluye \\
\midrule
Meta de plantación & El polígono debe alcanzar 658 plantas/ha contando existentes y sobrevivientes & Es la meta explícita del reto \\
Cantidad por especie & Cada especie debe quedar cerca de su proporción del plan, dentro del margen de ajuste & Sin ella el modelo podría vaciar una especie completa \\
Confiabilidad & Como la supervivencia es incierta, la meta solo puede garantizarse con cierta probabilidad & Comprar para el caso promedio incumpliría la meta con frecuencia \\
Presupuesto & El costo de la compra no debe superar el presupuesto del polígono & Límite económico explícito del reto \\
Existencia previa & Las plantas existentes cuentan hacia la meta & Evita comprar de más donde la meta ya se cumple en parte \\
Capacidad de la retícula & El número de posiciones del polígono es fijo, según su superficie y la densidad & No caben más plantas de las que admite el área \\
Compatibilidad entre especies & Ciertos pares de especies no deberían quedar como vecinos inmediatos & Es el criterio ecológico de la distribución espacial \\
\bottomrule
\end{tabularx}
\end{table*}

\section{Metodología de solución}
\label{sec:metodologia}

El problema se descompone en dos niveles que se resuelven en orden. El primero decide cuántas plantas de cada especie comprar en cada polígono; el segundo acomoda esas plantas en la retícula al tresbolillo. La \cref{subsec:formulacion} formula ambos niveles, la \cref{subsec:algoritmos} describe los algoritmos con que se resuelven y la \cref{sec:implementacion} resume su implementación.

\subsection{Formulación matemática}
\label{subsec:formulacion}

\subsubsection{Nivel 1: compra bajo incertidumbre}
\label{sec:modelo-compra}

El problema se descompone en dos niveles: la cantidad de plantas de cada especie que se adquiere en cada polígono y su acomodo en la retícula de plantación. Aquí se formula el primero, cuyo resultado es el dato de entrada del segundo (\cref{sec:modelo-espacial}). La compra se decide sin saber cuántas plantas nuevas sobrevivirán, por lo que se plantea como un programa estocástico con restricciones de probabilidad.

\paragraph{Datos y parámetros}
\label{subsec:compra-parametros}

Sea $\mathcal{S}$ el conjunto de las diez especies e $i$ un polígono de superficie $A_i$ (ha). Para cada $s\in\mathcal{S}$, $n_s^0$ es la densidad del plan de reforestación ($\sum_sn_s^0=658$ plantas/ha), $c_s$ el costo unitario de la planta sin logística y $q_s$ la supervivencia del primer año, observada en el monitoreo. $E_{is}$ es el número de plantas existentes de la especie $s$ en el polígono $i$ según los conteos de INECOL; en conjunto equivalen a 128.8 plantas/ha, 19.6\,\% de la meta, sobre los 183.01~ha de los 30 polígonos. La \cref{tab:compra-parametros} reporta estos valores. Se consideran además cuatro parámetros de política: el presupuesto por hectárea $B$, fijado por el socio formador en \$12{,}500; el sobrecosto logístico $\theta\in[0.05,0.10]$, como fracción del costo de la planta; la desviación máxima por especie respecto al plan, $\delta=0.10$; y el riesgo admisible de incumplimiento, $\alpha=0.05$.

\begin{table*}[t]
\centering
\caption{Parámetros por especie. $\bar{E}_s=\sum_iE_{is}/\sum_iA_i$: densidad media de plantas existentes; $c_s/q_s$: costo por planta sobreviviente.}
\label{tab:compra-parametros}
\small
\begin{tabular}{@{}lrrrrr@{}}
\toprule
Especie & $n_s^0$ (plantas/ha) & $c_s$ (\$) & $q_s$ & $\bar{E}_s$ (plantas/ha) & $c_s/q_s$ (\$) \\
\midrule
\textit{Agave lechuguilla}        &  42 & 26.0 & 0.943 &  8.1 & 27.6 \\
\textit{Agave salmiana}           & 196 & 26.0 & 0.959 & 38.0 & 27.1 \\
\textit{Agave scabra}             &  42 & 26.0 & 0.825 &  8.0 & 31.5 \\
\textit{Agave striata}            &  42 & 26.0 & 0.952 &  7.7 & 27.3 \\
\textit{Opuntia cantabrigiensis}  &  49 & 18.0 & 0.927 & 10.0 & 19.4 \\
\textit{Opuntia engelmannii}      &  38 & 21.0 & 0.899 &  8.0 & 23.4 \\
\textit{Opuntia robusta}          &  73 & 18.0 & 0.958 & 15.0 & 18.8 \\
\textit{Opuntia streptacantha}    &  64 & 18.0 & 0.886 & 12.5 & 20.3 \\
\textit{Prosopis laevigata}       &  86 & 26.5 & 0.821 & 16.5 & 32.3 \\
\textit{Yucca filifera}           &  26 & 26.0 & 0.884 &  5.0 & 29.4 \\
\midrule
Total                             & 658 &      &       & 128.8 &     \\
\bottomrule
\end{tabular}
\end{table*}

\paragraph{Formulación}
\label{subsec:compra-formulacion}

Sean $n_{is}\in\mathbb{Z}_{\geq0}$ la meta de plantas de la especie $s$ en el polígono $i$ y $x_{is}\in\mathbb{Z}_{\geq0}$ la cantidad a adquirir. La única variable aleatoria es el número de sobrevivientes de lo plantado, $S_{is}(x)\sim\operatorname{Bin}(x,q_s)$ (\cref{eq:binomial}), bajo independencia entre individuos. La \cref{tab:compra-notacion} resume la notación. Para cada polígono se resuelve
\begin{align}
\min_{n,\,x}\quad & (1+\theta)\sum_{s}c_s\Bigl(x_{is}+\mathbb{E}\!\left[D_{is}\right]\Bigr) \label{eq:obj}\\
\text{s.a.}\quad & \sum_{s}n_{is}=N_i, \label{eq:total}\\
& \left|n_{is}-n_s^0A_i\right|\leq\delta\,n_s^0A_i, && s\in\mathcal{S}, \label{eq:banda}\\
& (1+\theta)\sum_{s}c_sx_{is}\leq B\,A_i, \label{eq:presupuesto}\\
& \mathbb{P}\!\left(E_{is}+S_{is}(x_{is})\geq n_{is}\right)\geq1-\alpha, && s\in\mathcal{S}, \label{eq:probabilidad}
\end{align}
donde $N_i$ es $658A_i$ redondeado al entero más cercano y $D_{is}=\left(n_{is}-E_{is}-S_{is}(x_{is})\right)^{+}$ es el déficit al término del primer año. El objetivo \eqref{eq:obj} suma el costo de adquisición y el costo esperado de reponer el déficit, que representa la obligación contractual de reponer la planta muerta; la reposición se valora al mismo costo que la compra, logística incluida. La restricción \eqref{eq:total} fija la densidad total; \eqref{eq:banda} acota la composición por especie; \eqref{eq:presupuesto} limita el gasto, y \eqref{eq:probabilidad} exige que la meta de cada especie se alcance con probabilidad de al menos $1-\alpha$ \autocite{charnes1959}.

\begin{table*}[t]
\centering
\caption{Notación del modelo de compra.}
\label{tab:compra-notacion}
\small
\begin{tabularx}{\textwidth}{@{}l>{\RaggedRight\arraybackslash}p{0.17\textwidth}>{\RaggedRight\arraybackslash}X>{\RaggedRight\arraybackslash}p{0.2\textwidth}@{}}
\toprule
Símbolo & Tipo & Descripción & Unidad o valor \\
\midrule
$\mathcal{S}$, $s$ & Conjunto e índice & Especies a plantar & 10 especies \\
$i$ & Índice & Polígono de reforestación & 30 polígonos \\
$A_i$ & Dato & Superficie del polígono $i$ & ha \\
$N_i$ & Dato & Meta total del polígono, $658A_i$ redondeado & plantas \\
$n_s^0$ & Dato & Plantas de la especie $s$ por hectárea según el plan & plantas/ha \\
$c_s$ & Dato & Costo unitario de la planta, sin logística & \$/planta \\
$E_{is}$ & Dato & Plantas existentes de la especie $s$ en el polígono $i$ & plantas \\
$q_s$ & Parámetro estimado & Probabilidad de que una planta nueva sobreviva el primer año & $[0,1]$ \\
$B$ & Parámetro de política & Presupuesto por hectárea & \$12{,}500 \\
$\theta$ & Parámetro de política & Sobrecosto logístico como fracción del costo de la planta & $[0.05,\,0.10]$ \\
$\delta$ & Parámetro de política & Desviación relativa máxima de la meta de cada especie & 0.10 \\
$\alpha$ & Parámetro de política & Probabilidad admisible de no alcanzar la meta de una especie & 0.05 \\
$S_{is}(x)$ & Variable aleatoria & Sobrevivientes al primer año de $x$ plantas nuevas; $\operatorname{Bin}(x,q_s)$ & plantas \\
$D_{is}$ & Variable aleatoria & Déficit al término del primer año, $(n_{is}-E_{is}-S_{is})^{+}$ & plantas \\
$n_{is}$ & Variable de decisión & Meta de plantas de la especie $s$ en el polígono $i$ & plantas; entera \\
$x_{is}$ & Variable de decisión & Plantas de la especie $s$ a adquirir y plantar en el polígono $i$ & plantas; entera \\
\bottomrule
\end{tabularx}
\end{table*}

Antes de resolver conviene estimar el orden de magnitud. Con la composición del plan, el costo de cubrir el faltante, $\sum_sc_s(n_s^0-\bar{E}_s)$, asciende a \$12{,}456/ha, y a \$13{,}683/ha al corregirlo por la mortalidad del primer año, $\sum_sc_s(n_s^0-\bar{E}_s)/q_s$, aun sin margen de seguridad ni logística. Con $B=\$12{,}500$, la restricción \eqref{eq:presupuesto} es por tanto prácticamente activa, y la factibilidad depende de la banda \eqref{eq:banda}, que permite desplazar la composición hacia especies con menor costo por planta sobreviviente (\$18.8 para \textit{O.~robusta} frente a \$32.3 para \textit{P.~laevigata}). Cuando el problema es infactible, el modelo entrega el presupuesto mínimo $B_i^{*}$ con el que dejaría de serlo.

\paragraph{Supuestos y extensiones}
\label{subsec:compra-supuestos}

El modelo descansa en cuatro supuestos. (i) La supervivencia es independiente entre individuos; los eventos que afectan a toda la plantación, como la sequía, aumentan la varianza de $S_{is}$ y pueden incorporarse sorteando por escenario una tasa común $q_s^{(k)}$, aunque un solo año de monitoreo no permite estimar su dispersión. (ii) Los conteos de plantas existentes son exactos y esas plantas no mueren en el horizonte. Si el conteo fuera incierto, $E_{is}$ puede sortearse en cada escenario como $\operatorname{Poisson}(\bar{E}_sA_i)$: el estadístico de dispersión de Pearson de los conteos respecto a $\bar{E}_sA_i$ está entre 0.58 y 1.45 para todas las especies, sin evidencia de la sobredispersión que justificaría una binomial negativa \autocite{cameron2013}. (iii) El horizonte es de un año, aunque el monitoreo muestra que la mortalidad continúa el segundo; para obligaciones de reposición multianuales, el término de reposición se extiende con la cadena de Markov por posición de la \cref{subsec:markov}. (iv) El costo logístico es proporcional al costo de la planta; los costos fijos por viaje o por vivero se incorporan con variables binarias de activación sin alterar la estructura del modelo.

\section{Modelo de distribución espacial}
\label{sec:modelo-espacial}

El segundo nivel acomoda en el terreno las plantas adquiridas según la \cref{sec:modelo-compra} de modo que se minimice la competencia entre plantas vecinas. Se adopta la sobreplantación: se plantan todas las plantas compradas, y la mortalidad del primer año reduce la densidad hacia la meta. Siguiendo la \cref{subsec:coloracion}, el terreno es un grafo cuyos vértices son las posiciones de plantación y cuyas aristas unen posiciones vecinas, y el problema consiste en asignar una especie a cada vértice respetando la cantidad de cada especie \autocite{elizalde2025,eisenblatter2002}.

\subsection{Mosaico y objetivo}
\label{subsec:espacial-mosaico}

\paragraph{Mosaico periódico} Optimizar posición por posición un polígono de varios miles de plantas no es práctico ni se puede llevar a campo. En su lugar se optimiza un mosaico de pocas decenas de posiciones que la cuadrilla repite sobre todo el polígono como una tabla de claves de especie. El mosaico es un fragmento del tresbolillo de $4\times13=52$ posiciones con bordes periódicos: la última fila es vecina de la primera y la última columna de la primera, de modo que el mosaico se repite sin costuras y cada posición tiene exactamente seis vecinas, todas a la misma distancia $d\approx4.19$~m. El grafo tiene $52\cdot6/2=156$ aristas. La composición del mosaico, $m_s$ posiciones de la especie $s$ con $\sum_sm_s=52$, es proporcional a la compra $x_{is}$ del polígono, redondeada por el método del mayor residuo; los polígonos con la misma composición comparten mosaico.

\paragraph{Objetivo} La competencia total de una asignación $\sigma:V\to\mathcal{S}$ es
\begin{equation}
  C(\sigma)=\sum_{(u,v)\in\mathcal{E}}W_{\sigma(u)\sigma(v)},
  \label{eq:competencia}
\end{equation}
con $W$ la matriz de competencia de la \cref{subsec:espacial-matriz}. Como todas las aristas del tresbolillo representan la misma distancia, no se ponderan; un arreglo en cuadros, con vecinas diagonales más lejanas que las ortogonales, se trataría con un peso $w_{uv}$ por arista. Para comparar mosaicos de distinta composición se reporta el coeficiente de competencia $C(\sigma)/\mathbb{E}[C]$, donde
\begin{equation}
  \mathbb{E}[C]=\frac{|\mathcal{E}|}{n(n-1)}\sum_{a,b\in\mathcal{S}}m_a\bigl(m_b-\mathbf{1}[a=b]\bigr)W_{ab}
  \label{eq:competencia-azar}
\end{equation}
es la competencia esperada de una asignación aleatoria con la misma composición, con $n=52$: cada arista une dos posiciones distintas elegidas al azar, y el par de especies $(a,b)$ aparece con probabilidad $m_a(m_b-\mathbf{1}[a=b])/[n(n-1)]$. Un coeficiente menor que uno indica una reducción respecto al azar. El problema es una coloración de vértices con penalización y conteo fijo por color, NP-difícil \autocite{eisenblatter2002}, y se resuelve con dos métodos que comparten el objetivo \eqref{eq:competencia}: un algoritmo memético (\cref{subsec:espacial-memetico}) y un modelo de programación entera (\cref{subsec:espacial-entera}) cuya cota inferior mide la brecha del primero.

\subsection{Matriz de competencia}
\label{subsec:espacial-matriz}

Una matriz simétrica $W\in\mathbb{R}^{|\mathcal{S}|\times|\mathcal{S}|}$ cuantifica, para cada par de especies, la penalización por ocupar posiciones vecinas; su diagonal representa la competencia intraespecífica. Las funciones de correlación de pares, que miden asociación o segregación entre especies \autocite{wiegand2014}, requieren las coordenadas de cada planta, y los datos disponibles solo registran el número de plantas por especie en cada polígono. Como comprobación, las correlaciones entre los residuales de Pearson de los conteos de cada par de especies a lo largo de los 30 polígonos tienen media $-0.07$ y solo 4 de 45 son significativas al 5\,\%, cerca de lo esperado por azar, por lo que los conteos no contienen información sobre la interacción entre vecinas. La matriz se construye, en cambio, a partir de rasgos de las especies:
\begin{equation}
W_{ab}=\gamma^{\,\mathbf{1}[a=b]}\,T_{ab}\,N_{ab}-\phi\,F_{ab}.
\label{eq:matriz-W}
\end{equation}

El término de tamaño $T_{ab}=\min\{1,(r_a+r_b)/d\}$ es la fracción de la separación entre posiciones que cubren las copas adultas de las dos plantas, con $r_s$ el radio de copa de la especie $s$; vale 1 cuando las copas se tocan, de modo que dos especies grandes compiten más que dos pequeñas. El término de nicho $N_{ab}$ es el índice de traslape de \textcite{pianka1973} entre los perfiles de raíz por profundidad, $p_s=(p_{s1},p_{s2},p_{s3})$ con estratos de 0--30~cm, 30--100~cm y más de 100~cm,
\begin{equation}
N_{ab}=\frac{\sum_k p_{ak}\,p_{bk}}{\sqrt{\sum_k p_{ak}^2\,\sum_k p_{bk}^2}},
\end{equation}
que vale 1 cuando dos especies toman el agua del mismo estrato. Agaves y nopales concentran sus raíces en los primeros centímetros del suelo \autocite{snyman2006}, mientras que \textit{P.~laevigata} combina raíces superficiales con una raíz pivotante profunda, por lo que su traslape con las suculentas es menor. El factor $\gamma=1.5$ penaliza adicionalmente la vecindad entre individuos de la misma especie, que usan exactamente los mismos recursos y comparten plagas; que la competencia intraespecífica supere a la interespecífica es la condición para la coexistencia de las especies \autocite{chesson2000}, aunque su magnitud se fija por supuesto. Por último, $F_{ab}=1$ cuando uno de los dos individuos es \textit{P.~laevigata} y el otro una suculenta, y $\phi=0.3$ resta competencia por el efecto nodriza del mezquite, que favorece el establecimiento de otras especies bajo su copa \autocite{flores2003,padilla2006} y eleva la fertilidad del suelo \autocite{garciasanchez2012}.

La \cref{tab:matriz-competencia} presenta la matriz resultante, con valores entre 0.17 y 1.50. Los cinco valores más altos son intraespecíficos y corresponden a las especies de mayor tamaño (\textit{P.~laevigata}, \textit{O.~streptacantha}, \textit{Y.~filifera}, \textit{A.~salmiana} y \textit{O.~engelmannii}), mientras que \textit{P.~laevigata} tiene los valores más bajos con el resto de las especies, por la diferencia de raíces y el efecto nodriza. Así, el modelo prioriza separar a los individuos grandes de la misma especie y colocar las suculentas junto al mezquite. Como cada arista aporta exactamente un término al objetivo y el número de aristas es fijo, la asignación óptima no cambia ante una transformación afín positiva de la matriz, de modo que su escala no importa, aunque sí las diferencias relativas entre sus entradas. Los diámetros de copa y los perfiles de raíz son aproximados, y $\gamma$ y $\phi$ son supuestos; la sensibilidad de las claves a estos valores se mide en la \cref{subsec:res-sensibilidad}.

\begin{table}[t]
\centering
\caption{Matriz de competencia $W_{ab}$. ALE: \textit{A.~lechuguilla}; ASA: \textit{A.~salmiana}; ASC: \textit{A.~scabra}; AST: \textit{A.~striata}; OCA: \textit{O.~cantabrigiensis}; OEN: \textit{O.~engelmannii}; ORO: \textit{O.~robusta}; OST: \textit{O.~streptacantha}; PLA: \textit{P.~laevigata}; YFI: \textit{Y.~filifera}. En negritas, los valores mayores que 1.}
\label{tab:matriz-competencia}
\scriptsize
\setlength{\tabcolsep}{2.4pt}
\begin{tabular}{@{}lrrrrrrrrrr@{}}
\toprule
 & ALE & ASA & ASC & AST & OCA & OEN & ORO & OST & PLA & YFI \\
\midrule
ALE & 0.20 & 0.50 & 0.30 & 0.21 & 0.30 & 0.42 & 0.24 & 0.54 & 0.20 & 0.49 \\
ASA & 0.50 & \textbf{1.29} & 0.67 & 0.57 & 0.67 & 0.79 & 0.61 & 0.90 & 0.20 & 0.82 \\
ASC & 0.30 & 0.67 & 0.72 & 0.38 & 0.48 & 0.60 & 0.42 & 0.71 & 0.20 & 0.65 \\
AST & 0.21 & 0.57 & 0.38 & 0.43 & 0.38 & 0.50 & 0.32 & 0.62 & 0.20 & 0.56 \\
OCA & 0.30 & 0.67 & 0.48 & 0.38 & 0.72 & 0.60 & 0.42 & 0.72 & 0.17 & 0.63 \\
OEN & 0.42 & 0.79 & 0.60 & 0.50 & 0.60 & \textbf{1.07} & 0.54 & 0.84 & 0.17 & 0.74 \\
ORO & 0.24 & 0.61 & 0.42 & 0.32 & 0.42 & 0.54 & 0.54 & 0.66 & 0.17 & 0.58 \\
OST & 0.54 & 0.90 & 0.71 & 0.62 & 0.72 & 0.84 & 0.66 & \textbf{1.43} & 0.17 & 0.84 \\
PLA & 0.20 & 0.20 & 0.20 & 0.20 & 0.17 & 0.17 & 0.17 & 0.17 & \textbf{1.50} & 0.45 \\
YFI & 0.49 & 0.82 & 0.65 & 0.56 & 0.63 & 0.74 & 0.58 & 0.84 & 0.45 & \textbf{1.43} \\
\bottomrule
\end{tabular}
\end{table}

\subsection{Algoritmo memético}
\label{subsec:espacial-memetico}

El algoritmo combina un algoritmo genético, que explora mediante recombinación, con búsqueda tabú, que intensifica alrededor de cada descendiente; esta hibridación supera a cada componente por separado en coloración de grafos \autocite{galinier1999}.

\paragraph{Representación e inicialización} Un individuo asigna una especie a cada posición del mosaico respetando la composición $m_s$; su aptitud es la competencia total \eqref{eq:competencia}. La población inicial, de $P=20$ individuos, se forma con permutaciones aleatorias de la composición, y cada individuo se mejora con búsqueda tabú.

\paragraph{Selección y cruce} Los padres se eligen por torneo binario. El cruce adapta el operador GPX \autocite{galinier1999}: alternando entre padres, se elige la especie con más posiciones libres en el padre en turno y el hijo las hereda hasta completar su cantidad; las posiciones restantes se asignan de forma voraz a la especie que menos incrementa la competencia. A diferencia de la coloración, las especies no son intercambiables, por lo que el operador procede por especie y conserva la composición por construcción. Como dos mosaicos desplazados son el mismo esquema en el terreno, el segundo padre se traslada antes del cruce a la posición de máxima coincidencia con el primero.

\paragraph{Búsqueda tabú} Cada hijo se mejora con búsqueda tabú \autocite{glover1989,hertz1987} sobre intercambios de especie entre dos posiciones, que conservan la composición y cuyo efecto se evalúa solo en las vecinas de ambas. Se aplica el mejor intercambio no prohibido y se prohíbe revertirlo durante un número aleatorio de 5 a 15 iteraciones, salvo que produzca una nueva mejor solución; la búsqueda termina tras 500 iteraciones sin mejora.

\paragraph{Reemplazo y paro} El hijo sustituye al peor individuo si lo mejora y no es una copia de otro. La distancia entre individuos es el número de posiciones distintas minimizado sobre las traslaciones del mosaico, de modo que un esquema desplazado cuenta como copia; esto preserva la diversidad de la población. El algoritmo termina tras 1{,}000 generaciones o 100 sin mejora. Para cada mosaico se ejecutan 30 corridas independientes y se reportan la mejor competencia, su media y su desviación estándar.

\subsection{Modelo de programación entera}
\label{subsec:espacial-entera}

Sea $G=(V,\mathcal{E})$ el grafo del mosaico con bordes periódicos. Con $y_{vs}\in\{0,1\}$, igual a 1 si la posición $v$ recibe la especie $s$, y $z_{uvab}\geq0$, que representa el producto $y_{ua}y_{vb}$, se resuelve
\begin{align}
\min\quad & \sum_{(u,v)\in\mathcal{E}}\sum_{a,b\in\mathcal{S}}W_{ab}\,z_{uvab} \label{eq:ent-obj}\\
\text{s.a.}\quad & \textstyle\sum_{s}y_{vs}=1, && v\in V, \label{eq:ent-una}\\
& \textstyle\sum_{v}y_{vs}=m_s, && s\in\mathcal{S}, \label{eq:ent-conteo}\\
& \textstyle\sum_{b}z_{uvab}=y_{ua}, && (u,v)\in\mathcal{E},\ a\in\mathcal{S}, \label{eq:ent-lin1}\\
& \textstyle\sum_{a}z_{uvab}=y_{vb}, && (u,v)\in\mathcal{E},\ b\in\mathcal{S}. \label{eq:ent-lin2}
\end{align}
Las restricciones \eqref{eq:ent-una} y \eqref{eq:ent-conteo} imponen una especie por posición y la composición; \eqref{eq:ent-lin1}--\eqref{eq:ent-lin2} linealizan el objetivo cuadrático, pues con $y$ binaria fuerzan $z_{uvab}=y_{ua}y_{vb}$ sin depender del signo de $W$. Como el mosaico es invariante bajo traslaciones, se fija la especie más abundante en una posición, lo que elimina soluciones equivalentes sin pérdida de optimalidad. Con 52 posiciones, el modelo tiene 520 variables binarias y $156\cdot10^{2}=15{,}600$ continuas. Se resuelve con un solver de programación entera mixta con límite de tiempo, partiendo de la mejor solución del memético: esa solución inicial no cambia el óptimo ni la cota inferior, pero da desde el inicio una cota superior con la cual podar el árbol de búsqueda. La cota inferior que reporta el solver acota la brecha del memético respecto al óptimo, y las dos asignaciones se comparan posición por posición después de alinearlas por traslación.


\subsection{Algoritmos}
\label{subsec:algoritmos}

% Algoritmos de los dos niveles: curvas de compra, mochila, memético y
% resolución del modelo entero. La formulación está en modelo-compra.tex y
% modelo-espacial.tex.

\subsubsection{Curvas de compra por simulación Monte Carlo}
\label{subsec:compra-solucion}

Las especies solo se acoplan mediante \eqref{eq:total} y \eqref{eq:presupuesto}, lo que permite resolver el problema en dos etapas: primero, para cada especie y cada meta posible, la compra mínima y su costo; después, la combinación de metas de menor costo.

Para una meta $n$, el faltante es $L=n-E_{is}$; si $L\leq0$, las existentes bastan y $x^{*}=0$. En otro caso, por la \cref{eq:dualidad}, la menor compra que satisface \eqref{eq:probabilidad} es el cuantil $1-\alpha$ del número de plantas $X_L$ que hay que plantar hasta lograr $L$ sobrevivientes, con $X_L$ binomial negativa. El cuantil se estima por simulación Monte Carlo \autocite{:ross2013}: en cada escenario $k=1,\dots,K$ se simula la secuencia de supervivencias de las plantas, una tras otra, y se registra $x_k(L)$, la posición de la $L$-ésima sobreviviente, para todas las $L$ a la vez. La compra mínima estimada $\hat{x}^{*}(n)$ es el $r$-ésimo menor de los $x_k(L)$, con $r=\lceil(1-\alpha)K\rceil$: con esa compra, al menos $r$ escenarios alcanzan la meta, y con una planta menos no. Con los mismos escenarios se estima el déficit esperado $\hat{D}(n)$ como el promedio, sobre los $K$ escenarios, del faltante que deja la compra $\hat{x}^{*}(n)$, que es cero en los escenarios que alcanzan la meta. Como una sola secuencia por escenario sirve para todas las metas y todos los valores de $\alpha$, las comparaciones entre metas usan números aleatorios comunes, lo que reduce su varianza. La misma curva puede calcularse de forma exacta con la cola de la distribución binomial, que sirve para validar la simulación.

\paragraph{Convergencia y número de escenarios} El procedimiento es una aproximación por promedio muestral del programa \eqref{eq:obj}--\eqref{eq:probabilidad}, cuya solución converge a la del problema original al crecer $K$ \autocite{kleywegt2002,luedtke2008}. Para que el intervalo de confianza de 95\,\% de una probabilidad estimada $\hat{p}$ cercana a $0.95$ tenga un semiancho de $\varepsilon=0.005$ se requiere $K\geq z_{0.025}^{2}p(1-p)/\varepsilon^{2}\approx7{,}300$ (\cref{eq:intervalo}); se adopta $K=10{,}000$.

\subsubsection{Mochila de elección múltiple}
\label{subsec:mochila}

Elegir una meta $n_s$ por especie, dentro de su banda \eqref{eq:banda}, que minimice $\sum_sC_s(n_s)$ con $C_s(n)=(1+\theta)c_s\bigl(\hat{x}_s^{*}(n)+\hat{D}_s(n)\bigr)$, sujeto a \eqref{eq:total} y \eqref{eq:presupuesto}, es un problema de mochila de elección múltiple. Sin la restricción de presupuesto se resuelve de forma exacta con programación dinámica, agregando una especie a la vez y llevando como estado la suma acumulada de metas; el costo es $O(|\mathcal{S}|\,H\,w)$, con $H$ la suma de los anchos de las bandas y $w$ el mayor de ellos. La misma recursión, minimizando solo el costo de compra, da el presupuesto mínimo $B_i^{*}$. Si el óptimo sin presupuesto cabe en $BA_i$, es la solución; si $B<B_i^{*}$, el polígono es infactible; solo cuando $B$ queda entre ambos valores se resuelve el problema completo como programa entero mixto.

\subsubsection{Algoritmo memético}
\label{subsec:espacial-memetico}

El algoritmo combina un algoritmo genético, que explora mediante recombinación, con búsqueda tabú, que intensifica alrededor de cada descendiente; esta hibridación supera a cada componente por separado en coloración de grafos \autocite{galinier1999}.

\paragraph{Representación e inicialización} Un individuo asigna una especie a cada posición del mosaico respetando la composición $m_s$; su aptitud es la competencia total \eqref{eq:competencia}. La población inicial, de $P=20$ individuos, se forma con permutaciones aleatorias de la composición, y cada individuo se mejora con búsqueda tabú.

\paragraph{Selección y cruce} Los padres se eligen por torneo binario. El cruce adapta el operador GPX \autocite{galinier1999}: alternando entre padres, se elige la especie con más posiciones libres en el padre en turno y el hijo las hereda hasta completar su cantidad; las posiciones restantes se asignan de forma voraz a la especie que menos incrementa la competencia. A diferencia de la coloración, las especies no son intercambiables, por lo que el operador procede por especie y conserva la composición por construcción. Como dos mosaicos desplazados son el mismo esquema en el terreno, el segundo padre se traslada antes del cruce a la posición de máxima coincidencia con el primero.

\paragraph{Búsqueda tabú} Cada hijo se mejora con búsqueda tabú \autocite{glover1989,hertz1987} sobre intercambios de especie entre dos posiciones, que conservan la composición y cuyo efecto se evalúa solo en las vecinas de ambas. Se aplica el mejor intercambio no prohibido y se prohíbe revertirlo durante un número aleatorio de 5 a 15 iteraciones, salvo que produzca una nueva mejor solución; la búsqueda termina tras 500 iteraciones sin mejora.

\paragraph{Reemplazo y paro} El hijo sustituye al peor individuo si lo mejora y no es una copia de otro. La distancia entre individuos es el número de posiciones distintas minimizado sobre las traslaciones del mosaico, de modo que un esquema desplazado cuenta como copia; esto preserva la diversidad de la población. El algoritmo termina tras 1{,}000 generaciones o 100 sin mejora. Para cada mosaico se ejecutan 30 corridas independientes y se reportan la mejor competencia, su media y su desviación estándar.

\subsubsection{Resolución del modelo entero}
\label{subsec:ple-solucion}

El modelo entero del mosaico (\cref{eq:ent-obj,eq:ent-una,eq:ent-conteo,eq:ent-lin1,eq:ent-lin2}) se resuelve con HiGHS. Se resuelve con un solver de programación entera mixta con límite de tiempo, partiendo de la mejor solución del memético: esa solución inicial no cambia el óptimo ni la cota inferior, pero da desde el inicio una cota superior con la cual podar el árbol de búsqueda. La cota inferior que reporta el solver acota la brecha del memético respecto al óptimo, y las dos asignaciones se comparan posición por posición después de alinearlas por traslación.


\section{Implementación}
\label{sec:implementacion}

Los dos modelos se implementaron en C++20 como una biblioteca, \texttt{refor\_core}, sobre la que se construyen tres programas: \texttt{refor}, una herramienta de línea de comandos para análisis y reportes; \texttt{refor\_ple}, que resuelve por separado el modelo entero de la \cref{subsec:espacial-entera}; y \texttt{refor\_gui}, una aplicación de escritorio para los técnicos de planeación. Los programas enteros se resuelven con HiGHS. El código, los datos procesados y las instrucciones para reproducir cada resultado de este trabajo están en el repositorio del proyecto.

\subsection{Organización}
\label{subsec:impl-organizacion}

La \cref{tab:modulos} relaciona cada módulo de la biblioteca con la parte del modelo que implementa. El flujo sigue la descomposición de las \cref{sec:modelo-compra,sec:modelo-espacial}: se leen los datos de INECOL, se calculan las curvas $\hat{x}_s^{*}(n)$ y $\hat{D}_s(n)$ de cada especie y polígono, se resuelve la mochila de cada polígono, se reparte su compra en un mosaico de 52 posiciones y se agrupan los polígonos con la misma composición, de modo que cada mosaico distinto se optimiza una sola vez.

\begin{table}[t]
\centering
\caption{Módulos de \texttt{refor\_core} y la parte del modelo que implementan.}
\label{tab:modulos}
\small
\begin{tabularx}{\columnwidth}{@{}>{\RaggedRight\arraybackslash}p{2.6cm}>{\RaggedRight\arraybackslash}X@{}}
\toprule
Módulo & Contenido \\
\midrule
\texttt{simulacion} & Escenarios Monte Carlo y curvas $\hat{x}^{*}(n)$, $\hat{D}(n)$ \\
\texttt{exacto} & Las mismas curvas con la cola binomial exacta \\
\texttt{mochila} & Programación dinámica y programa entero de la mochila de elección múltiple \\
\texttt{modelo} & Curvas en paralelo y solución de cada polígono \\
\texttt{plan} & Composición por mayor residuo, agrupación de polígonos y corridas del memético \\
\texttt{memetico/} & Mosaico toroidal, matriz $W$, búsqueda tabú y algoritmo memético \\
\texttt{asignacion\_exacta/} & Modelo entero del mosaico y comparación con el memético \\
\bottomrule
\end{tabularx}
\end{table}

\subsection{Reproducibilidad}
\label{subsec:impl-reproducibilidad}

Todo resultado depende de una sola semilla (42 por defecto) y debe ser idéntico sin importar el número de hilos ni la plataforma. Para lograrlo, cada flujo aleatorio se inicializa a partir de su identidad lógica y no del orden en que se ejecuta: los escenarios de la especie $s$ en el polígono $i$ usan una semilla derivada de $(42,i,s)$, y cada corrida del memético, una derivada del mosaico y del número de corrida. El generador es \texttt{mt19937\_64}, cuya secuencia fija el estándar de C++. En cambio, los algoritmos de las distribuciones de la biblioteca estándar difieren entre implementaciones, por lo que los enteros uniformes y las permutaciones se generan con funciones propias, y la supervivencia de cada planta se decide comparando la salida del generador con $q_s\cdot2^{64}$. La biblioteca se compila sin contracción de operaciones de punto flotante, que podría cambiar el redondeo entre compiladores. Como comprobación, la compra de los 30 polígonos calculada con un hilo y con diez da archivos idénticos byte por byte. El flujo de publicación compila los programas en Linux, macOS y Windows y compara sus resultados.

\subsection{Pruebas y desempeño}
\label{subsec:impl-pruebas}

La biblioteca tiene 30 casos de prueba automatizados. Los más importantes comparan cada método con una referencia independiente: la cola binomial con valores conocidos; las curvas de Monte Carlo con las exactas; la programación dinámica y el programa entero de la mochila con la enumeración exhaustiva de metas, con y sin presupuesto; y el modelo entero del mosaico con la enumeración exhaustiva en un mosaico de $3\times3$. Otras pruebas verifican invariantes: la monotonía de $\hat{x}^{*}(n)$ en $n$ y en $\alpha$, que la búsqueda tabú nunca empeora su solución inicial, que el cruce conserva la composición, que la distancia entre mosaicos reconoce un mismo esquema desplazado y que la secuencia aleatoria del memético es la misma en todas las plataformas.

En una computadora con un procesador Apple M4 de diez núcleos, las curvas de los 30 polígonos con $K=10{,}000$ escenarios se calculan en 1.8~s, y en 0.15~s con la cola binomial exacta. Las 30 corridas del memético para cada uno de los seis mosaicos distintos toman en conjunto unos 30~s. El modelo entero es el único paso costoso: se le dan 300~s por mosaico.


\section{Experimentación y resultados}
\label{sec:experimentacion}

Esta sección describe el equipo y el tamaño del problema, reporta la simulación del proceso de siembra en una hectárea, valida los dos métodos aproximados contra los exactos y presenta los resultados de los 30 polígonos. Todos los resultados se reproducen con \texttt{refor} y los scripts del repositorio del proyecto (\cref{sec:implementacion}), con la semilla 42.

\subsection{Características del equipo y del software}
\label{subsec:equipo}

Todos los experimentos se corrieron en el equipo de la \cref{tab:equipo}. Los programas se compilaron en modo de optimización (\texttt{Release}) y usan todos los núcleos disponibles; los resultados no dependen del número de hilos (\cref{subsec:impl-reproducibilidad}).

\begin{table}[t]
\centering
\caption{Equipo y software utilizados.}
\label{tab:equipo}
\small
\begin{tabularx}{\columnwidth}{@{}>{\RaggedRight\arraybackslash}p{2.5cm}>{\RaggedRight\arraybackslash}X@{}}
\toprule
Característica & Valor \\
\midrule
Modelo & MacBook Pro (Mac16,1) \\
Sistema operativo & macOS 26.5.2 \\
Disco duro & 494~GB, de estado sólido \\
Memoria RAM & 16~GB \\
Procesador & Apple M4 (ARM64) \\
Núcleos & 10 (4 de rendimiento y 6 de eficiencia) \\
\midrule
Lenguaje y compilador & C++20, Apple clang 17.0.0 \\
Construcción & CMake 4.3.0 \\
Solver entero & HiGHS 1.14.0 \\
Bibliotecas & CLI11 2.4.2, Catch2 3.7.1; Dear ImGui 1.92.8, ImPlot 1.0 y GLFW 3.4 (interfaz) \\
Análisis & Python 3.14.8 \\
Reporte & TeX Live 2025, pdf\TeX{} 1.40.27, Biber 2.20 \\
\bottomrule
\end{tabularx}
\end{table}

\subsection{Tamaño del problema}
\label{subsec:tamano}

La \cref{tab:tamano} resume el tamaño de la instancia. Hay 10 especies y una meta de 658 plantas por hectárea, que en los 30 polígonos suma 120{,}419 plantas; la composición del plan por especie está en la \cref{tab:compra-parametros} y la que resulta de la compra, en la \cref{tab:siembra}. En el nivel de compra, cada polígono tiene 20 variables enteras (una meta y una compra por especie); al resolverlo como mochila de elección múltiple hay una variable binaria por cada meta posible de cada especie, 24{,}094 en total, y las curvas que alimentan la mochila se estiman con 10{,}000 escenarios para cada uno de los 300 pares de polígono y especie. En el nivel espacial, el memético decide la especie de 52 posiciones por mosaico, y el modelo entero de cada mosaico tiene 520 variables binarias, 15{,}600 continuas y 3{,}182 restricciones.

\begin{table}[t]
\centering
\caption{Tamaño del problema. Los parámetros del modelo son datos o supuestos; los de los algoritmos controlan su ejecución.}
\label{tab:tamano}
\small
\begin{tabularx}{\columnwidth}{@{}>{\RaggedRight\arraybackslash}X r@{}}
\toprule
Elemento & Cantidad \\
\midrule
Especies & 10 \\
Polígonos (superficie total) & 30 (183.01~ha) \\
Plantas por hectárea (meta) & 658 \\
Meta total en los 30 polígonos & 120{,}419 \\
Plantas existentes registradas & 23{,}569 \\
\midrule
\multicolumn{2}{@{}l}{\textit{Variables}} \\
Compra: metas y compras por polígono & 20 \\
Compra: total en los 30 polígonos & 600 \\
Mochila: binarias en los 30 polígonos & 24{,}094 \\
Escenarios Monte Carlo por curva & 10{,}000 \\
Espacial: posiciones por mosaico & 52 \\
Modelo entero: binarias / continuas & 520 / 15{,}600 \\
Modelo entero: restricciones & 3{,}182 \\
\midrule
\multicolumn{2}{@{}l}{\textit{Parámetros del modelo}} \\
Por especie ($n_s^0$, $c_s$, $q_s$) & 30 \\
Por polígono ($A_i$) & 30 \\
Existentes por polígono y especie ($E_{is}$) & 300 \\
Matriz de competencia (entradas distintas) & 55 \\
De política ($B$, $\theta$, $\delta$, $\alpha$) & 4 \\
Total & 419 \\
\midrule
\multicolumn{2}{@{}l}{\textit{Parámetros de los algoritmos}} \\
Escenarios, semilla & 2 \\
Memético: población, generaciones, paro, corridas & 4 \\
Tabú: iteraciones sin mejora, permanencia mínima y máxima & 3 \\
Mosaico: filas y columnas & 2 \\
Modelo entero: límite de tiempo & 1 \\
\bottomrule
\end{tabularx}
\end{table}

\subsection{Simulación del proceso de siembra}
\label{subsec:simulacion}

Para estimar qué ocurre en una hectárea típica se simuló 1{,}000 veces el proceso de siembra con el script \texttt{simulacion\_siembra.py}, que lee los datos de INECOL y las salidas de \texttt{refor}. La hectárea es representativa del sitio: su meta por especie, $n_s$, es la composición óptima de la compra sumada sobre los 30 polígonos y dividida entre su superficie, redondeada por mayor residuo a 658 plantas. En cada réplica:
\begin{enumerate}
  \item las plantas existentes de cada especie se sortean como $E_s\sim\operatorname{Poisson}(\bar{E}_s)$, el modelo de conteo que respalda el análisis de dispersión de la \cref{sec:modelo-compra};
  \item las plantas a suministrar son las que hay que plantar, una tras otra, hasta que sobrevivan $n_s-E_s$, una variable binomial negativa (\cref{eq:binneg});
  \item se fija la compra $x_s$ como el cuantil 95\,\% de las plantas a suministrar, la misma regla del modelo de compra con $\alpha=0.05$ pero con las existentes inciertas, y se sortea de nuevo la supervivencia, $S_s\sim\operatorname{Bin}(x_s,q_s)$, de modo que las plantas vivas al año son $E_s+S_s$;
  \item una hectárea al tresbolillo de $14\times47=658$ posiciones, con 1{,}853 pares de vecinas, se llena con una clave, cada planta sobrevive con probabilidad $q_s$, y la competencia es la suma de $W$ sobre los pares de vecinas que siguen vivas.
\end{enumerate}

\paragraph{Plantas existentes, a suministrar y vivas} La \cref{tab:siembra} reporta los valores esperados por especie. El valor esperado del total de plantas existentes es de 128.6 por hectárea, con desviación estándar de 11.4, igual al promedio observado (128.8) como se espera del modelo Poisson. Para alcanzar exactamente la meta hay que suministrar en promedio 579.1 plantas por hectárea, pero con ese promedio cada especie cumple cerca de la mitad de las veces; para cumplir con probabilidad de 95\,\% en cada especie hay que suministrar 651, y con ellas la hectárea tiene en promedio 722.4 plantas vivas al año. Cada especie cumple su meta en 95 a 98\,\% de las réplicas, pero las diez a la vez solo en 68.4\,\%: la restricción \eqref{eq:probabilidad} garantiza el cumplimiento de cada especie por separado, no de todas juntas. Las 651 plantas de esta hectárea aislada superan las 594 por hectárea que compra \texttt{refor} en los 30 polígonos porque el margen de seguridad crece con la raíz de la meta y pesa más en una superficie chica, el mismo efecto que encarece los polígonos pequeños (\cref{subsec:res-compra}).

\begin{table*}[t]
\centering
\caption{Simulación de 1{,}000 siembras de una hectárea: meta por especie, valores esperados (desviación estándar) de las plantas existentes y de las plantas a suministrar para alcanzar la meta, compra que cumple la meta con probabilidad de 95\,\%, plantas vivas al año con esa compra y fracción de réplicas que cumplen la meta. En la fila total, la última columna es la fracción en que cumplen las diez especies a la vez.}
\label{tab:siembra}
\small
\begin{tabular}{@{}lrrrrrr@{}}
\toprule
Especie & Meta & Existentes & A suministrar & Compra al 95\,\% & Vivas al año & Cumple \\
\midrule
\textit{Agave lechuguilla}       &  39 &   8.2 (2.9) &  32.6 (3.5) &  38 &  44.1 (3.2) & 0.966 \\
\textit{Agave salmiana}          & 196 &  38.0 (6.0) & 164.8 (6.7) & 175 & 205.8 (6.4) & 0.949 \\
\textit{Agave scabra}            &  38 &   8.0 (2.8) &  36.4 (4.3) &  44 &  44.2 (3.7) & 0.962 \\
\textit{Agave striata}           &  39 &   7.7 (2.8) &  32.9 (3.2) &  38 &  43.9 (3.0) & 0.968 \\
\textit{Opuntia cantabrigiensis} &  54 &  10.1 (3.3) &  47.4 (4.0) &  54 &  60.2 (3.7) & 0.966 \\
\textit{Opuntia engelmannii}     &  41 &   8.0 (2.9) &  36.8 (3.9) &  43 &  46.6 (3.6) & 0.955 \\
\textit{Opuntia robusta}         &  80 &  14.9 (3.8) &  67.9 (4.3) &  75 &  86.6 (4.2) & 0.963 \\
\textit{Opuntia streptacantha}   &  70 &  12.4 (3.6) &  64.9 (5.0) &  73 &  77.0 (4.4) & 0.962 \\
\textit{Prosopis laevigata}      &  77 &  16.5 (4.1) &  73.8 (6.4) &  84 &  85.3 (5.3) & 0.952 \\
\textit{Yucca filifera}          &  24 &   4.9 (2.2) &  21.6 (3.1) &  27 &  28.8 (2.6) & 0.980 \\
\midrule
Total                            & 658 & 128.6 (11.4) & 579.1 (14.5) & 651 & 722.4 (13.2) & 0.684 \\
\bottomrule
\end{tabular}
\end{table*}

\paragraph{Competencia en una hectárea} La \cref{tab:siembra-competencia} compara la competencia de la hectárea llena con cuatro arreglos: la clave 1 del memético (\cref{fig:clave1}), la secuencia vigente con las filas alineadas, la vigente con el mejor desfase entre filas (\cref{tab:clave-vigente}) y una asignación al azar con la composición de la clave 1. El valor esperado de la competencia al año con la clave 1 es de 802.7 por hectárea, 18.0\,\% menos que al azar (979.0). La secuencia vigente con filas alineadas llega a 1{,}090.4, 11.4\,\% más que al azar y 35.8\,\% más que la clave 1; con el mejor desfase baja a 842.7, todavía 5.0\,\% por encima de la clave 1. Los intervalos de 90\,\% de la clave 1 y de la vigente alineada no se traslapan, de modo que la diferencia no se debe al azar de la mortalidad. La mortalidad reduce la competencia en alrededor de 15\,\% respecto a la hectárea recién plantada, porque cada planta muerta elimina hasta seis pares de vecinas.

\begin{table}[t]
\centering
\caption{Competencia en una hectárea de 658 posiciones: al plantar y valor esperado al año (desviación estándar e intervalo de 90\,\%) en 1{,}000 réplicas de la mortalidad.}
\label{tab:siembra-competencia}
\footnotesize
\setlength{\tabcolsep}{3pt}
\begin{tabular}{@{}lrrr@{}}
\toprule
Arreglo & Al plantar & Al año & Int. 90\,\% \\
\midrule
Clave 1              &  945 &  803 (18) &   773--832 \\
Vigente, mejor desf. &  990 &  843 (20) &   809--874 \\
Al azar              & 1159 &  979 (25) &   939--1022 \\
Vigente, alineada    & 1304 & 1090 (27) &  1047--1136 \\
\bottomrule
\end{tabular}
\end{table}

\paragraph{Probabilidad de transición entre plantas} Si cada planta es un estado, su especie, el orden de plantación define una cadena de Markov (\cref{subsec:markov}): $P_{ij}$ es la probabilidad de que la siguiente planta de la fila sea de la especie $j$ cuando la actual es de la $i$. La \cref{tab:transicion} da la matriz de la clave 1, estimada con las 52 transiciones de sus cuatro filas, que se repiten a lo largo del surco. La diagonal es nula: ninguna planta tiene a una de su especie inmediatamente después. Varias transiciones son casi deterministas y reflejan la estructura de la clave: después de \textit{P.~laevigata} sigue \textit{O.~streptacantha} con probabilidad 0.71, después de \textit{O.~streptacantha} siempre sigue \textit{A.~salmiana}, y \textit{A.~salmiana} va seguida de \textit{P.~laevigata} con probabilidad 0.40. La secuencia vigente también tiene la diagonal nula, lo que muestra que la matriz de transición a lo largo de la fila no basta para juzgar un arreglo: la vigente alineada evita repetir especie dentro de la fila, pero la repite en las filas vecinas, que es donde concentra su competencia.

\begin{table*}[t]
\centering
\caption{Matriz de transición entre especies a lo largo del orden de plantación de la clave 1: probabilidad de que la planta siguiente sea de la especie de la columna dado que la actual es de la del renglón. Abreviaturas como en la \cref{tab:matriz-competencia}; los ceros se omiten.}
\label{tab:transicion}
\small
\begin{tabular}{@{}lrrrrrrrrrr@{}}
\toprule
 & ALE & ASA & ASC & AST & OCA & OEN & ORO & OST & PLA & YFI \\
\midrule
ALE &      &      &      &      &      &      & 1.00 &      &      &      \\
ASA &      &      & 0.07 & 0.20 & 0.13 &      & 0.07 &      & 0.40 & 0.13 \\
ASC & 0.67 & 0.33 &      &      &      &      &      &      &      &      \\
AST &      &      & 0.33 &      &      & 0.33 & 0.33 &      &      &      \\
OCA & 0.25 & 0.75 &      &      &      &      &      &      &      &      \\
OEN &      & 0.33 &      &      &      &      & 0.33 &      & 0.33 &      \\
ORO &      & 0.67 &      &      & 0.17 & 0.17 &      &      &      &      \\
OST &      & 1.00 &      &      &      &      &      &      &      &      \\
PLA &      &      &      &      & 0.14 & 0.14 &      & 0.71 &      &      \\
YFI &      &      & 0.50 &      &      &      &      & 0.50 &      &      \\
\bottomrule
\end{tabular}
\end{table*}

\paragraph{Orden de plantación} La \cref{tab:orden} da el orden en que se planta la clave 1, el ejemplo para los 18 polígonos que la comparten. Cada fila del terreno sigue la secuencia de 13 posiciones de su renglón, con 4.19~m entre plantas, y vuelve a empezar al terminarla, de modo que un surco de 100~m lleva 24 plantas: las 13 de la secuencia y las 11 primeras otra vez. Las filas se separan 3.63~m y cada una empieza 2.10~m, media separación, más adelante que la anterior, lo que forma el tresbolillo. Después de la fila 4 se repite la fila 1. La cuadrilla solo necesita la tabla de cuatro renglones; las vecinas de cada planta, en la fila y entre filas, quedan fijadas por ella (\cref{fig:clave1}).

\begin{table*}[t]
\centering
\caption{Orden de plantación de la clave 1. Cada renglón es una fila del terreno; la secuencia se repite a lo largo de la fila, y después de la fila 4 se repite la fila 1. Separación entre plantas de 4.19~m; entre filas, 3.63~m con un desfase de 2.10~m.}
\label{tab:orden}
\small
\setlength{\tabcolsep}{4pt}
\begin{tabular}{@{}lccccccccccccc@{}}
\toprule
Fila & 1 & 2 & 3 & 4 & 5 & 6 & 7 & 8 & 9 & 10 & 11 & 12 & 13 \\
\midrule
1 & Asa & Pla & Ost & Asa & Oca & Asa & Pla & Ost & Asa & Oro & Asa & Ast & Oro \\
2 & Oen & Asa & Asc & Asa & Pla & Ost & Asa & Yfi & Asc & Ale & Oro & Asa & Ast \\
3 & Oro & Oen & Pla & Ost & Asa & Pla & Oen & Oro & Oca & Asa & Ast & Asc & Ale \\
4 & Yfi & Ost & Asa & Pla & Ost & Asa & Oca & Asa & Pla & Oca & Ale & Oro & Asa \\
\bottomrule
\end{tabular}
\end{table*}

\subsection{Validación}
\label{sec:validacion}

Cada nivel se resolvió con un método aproximado (simulación Monte Carlo para la compra y un algoritmo memético para la distribución) y se validó contra un método exacto que comparte el mismo objetivo. Salvo que se indique otra cosa, los resultados usan $\alpha=0.05$, $\theta=0.05$, $\delta=0.10$, $K=10{,}000$ y la semilla 42.

\subsubsection{Compra: Monte Carlo contra la cola binomial}
\label{subsec:val-compra}

Como la única variable aleatoria es $S_{is}(x)\sim\operatorname{Bin}(x,q_s)$, la restricción \eqref{eq:probabilidad} puede evaluarse de forma exacta con la cola binomial, lo que permite medir el error de la simulación sin depender de otra simulación. Con ambos métodos se resolvió la compra de los 30 polígonos.

\paragraph{Curvas} De los 300 pares de polígono y especie, en 181 ambos métodos eligen la misma meta $n_{is}$, y en ellos la compra difiere a lo más en una planta (en 12 pares), lo esperado al estimar un cuantil con $10{,}000$ escenarios. En los 119 pares restantes, la mochila elige metas distintas entre soluciones de costo casi idéntico: la compra total es de $108{,}680$ plantas con Monte Carlo y de $108{,}679$ con el método exacto, y su costo difiere en \$13 sobre \$2.65~millones.

\paragraph{Presupuesto mínimo} El presupuesto mínimo ponderado por superficie, $\bar{B}^{*}=\sum_iA_iB_i^{*}/\sum_iA_i$, difiere entre ambos métodos a lo más en \$1.6/ha para los siete valores de $\alpha$ entre 0.01 y 0.50 (\cref{tab:frontera}), y en \$0.08/ha con $\alpha=0.05$ (\$14{,}496.10 frente a \$14{,}496.02).

\paragraph{Número de escenarios y semilla} Con $\alpha=0.05$, $\bar{B}^{*}$ vale \$14{,}493.90/ha con $K=1{,}000$, \$14{,}496.10/ha con $K=10{,}000$ y \$14{,}496.31/ha con $K=100{,}000$, frente a \$14{,}496.02/ha con el método exacto. Con $K=10{,}000$ y las semillas 1 a 5, $\bar{B}^{*}$ va de \$14{,}495.82 a \$14{,}496.83/ha. El error de Monte Carlo es, por tanto, del orden de \$1/ha, tres órdenes de magnitud por debajo de la diferencia entre el presupuesto disponible y el mínimo (\cref{subsec:res-compra}), y $K=10{,}000$ es suficiente.

\subsubsection{Distribución: memético contra el modelo entero}
\label{subsec:val-espacial}

Cada uno de los seis mosaicos distintos (\cref{subsec:res-espacial}) se resolvió también con el modelo entero de la \cref{subsec:espacial-entera}, con HiGHS, un límite de 300~s y la mejor clave del memético como solución inicial. La \cref{tab:val-espacial} reporta el resultado.

\begin{table}[t]
\centering
\caption{Memético contra el modelo entero en los seis mosaicos. Memético: mejor competencia de 30 corridas y su desviación estándar; cota: cota inferior del modelo entero a los 300~s; brecha: distancia máxima del memético al óptimo, $(C-\text{cota})/C$.}
\label{tab:val-espacial}
\small
\begin{tabular}{@{}crrrr@{}}
\toprule
Mosaico & Memético & d.e. & Cota & Brecha (\%) \\
\midrule
1 & 78.18 & 0.005 & 74.70 & 4.5 \\
2 & 77.37 & 0.006 & 74.27 & 4.0 \\
3 & 81.29 & 0.000 & 78.25 & 3.7 \\
4 & 76.77 & 0.005 & 73.55 & 4.2 \\
5 & 76.80 & 0.005 & 73.71 & 4.0 \\
6 & 81.85 & 0.000 & 78.82 & 3.7 \\
\bottomrule
\end{tabular}
\end{table}

En ningún mosaico el solver encontró una asignación mejor que la del memético, y en ninguno logró probar la optimalidad: la competencia del memético está a lo más entre 3.7\,\% y 4.5\,\% del óptimo. Esa brecha es una cota superior del error, y lo más probable es que el error real sea mucho menor. La cota inferior llega en los primeros 4~s a menos de 0.2 de su valor final y después casi no se mueve (en el mosaico 1 pasa de 74.69 a 74.70 en el resto de los 300~s), lo que indica que la brecha refleja la debilidad de la relajación lineal del modelo, no una mala solución. Es lo mismo que \textcite{eisenblatter2002} reportan para la asignación de frecuencias, donde las relajaciones lineales dan cotas débiles. Además, el memético es muy estable: en los seis mosaicos la desviación estándar de las 30 corridas no excede 0.006, y en los mosaicos 3 y 6 todas las corridas terminan en la misma competencia. Las corridas terminan, en promedio, entre la generación 100 y la 124, de modo que la última mejora ocurre antes de la generación 24 y el criterio de 100 generaciones sin mejora es holgado; en el mosaico 3, la búsqueda tabú sobre la población inicial ya encuentra la mejor clave en todas las corridas.

\subsection{Resultados en los 30 polígonos}
\label{sec:resultados}

Ambos niveles se resolvieron para los 30 polígonos (183.01~ha) con los parámetros de la \cref{sec:validacion}, el mosaico de $4\times13=52$ posiciones y la matriz de la \cref{tab:matriz-competencia}. Todos los resultados se reproducen con \texttt{refor} y la semilla 42.

\subsubsection{Compra}
\label{subsec:res-compra}

\paragraph{El presupuesto no alcanza} Con $B=\$12{,}500$/ha ningún polígono es factible. El presupuesto mínimo es, en promedio ponderado por superficie, $\bar{B}^{*}=\$14{,}496$/ha, 16\,\% más que el disponible: va de \$14{,}361/ha en el polígono 8 a \$14{,}934/ha en el polígono 1. Cumplir la meta en los 30 polígonos con $\alpha=0.05$ requiere \$365~mil adicionales. La infactibilidad es estructural, como anticipaba la estimación de la \cref{subsec:compra-formulacion}: el costo de cubrir el faltante corregido por mortalidad ya excedía el presupuesto antes de agregar margen de seguridad y logística. Aceptar más riesgo apenas ayuda (\cref{tab:frontera}): pasar de $\alpha=0.01$ a $\alpha=0.50$ ahorra 3.7\,\% del presupuesto a cambio de multiplicar por cincuenta la probabilidad de incumplir la meta de cada especie, y aun con $\alpha=0.50$ faltan \$296~mil. Lo que excede el presupuesto es el faltante mismo, no el margen de seguridad. El sobrecosto logístico pesa más que el riesgo: con $\theta=0.10$ en lugar de $0.05$, $\bar{B}^{*}$ sube a \$15{,}186/ha.

\begin{table}[t]
\centering
\caption{Presupuesto mínimo ponderado por superficie, $\bar{B}^{*}$ (\$/ha), con Monte Carlo y con la cola binomial exacta, y faltante respecto a \$12{,}500/ha en los 183.01~ha.}
\label{tab:frontera}
\small
\begin{tabular}{@{}rrrr@{}}
\toprule
$\alpha$ & Monte Carlo & Exacto & Faltante (miles de \$) \\
\midrule
0.01 & 14{,}664 & 14{,}666 & 396 \\
0.02 & 14{,}596 & 14{,}597 & 384 \\
0.05 & \textbf{14{,}496} & \textbf{14{,}496} & \textbf{365} \\
0.10 & 14{,}408 & 14{,}409 & 349 \\
0.20 & 14{,}305 & 14{,}305 & 330 \\
0.30 & 14{,}233 & 14{,}233 & 317 \\
0.50 & 14{,}117 & 14{,}118 & 296 \\
\bottomrule
\end{tabular}
\end{table}

\paragraph{Compra sin presupuesto} Sin la restricción \eqref{eq:presupuesto}, la compra óptima suma $108{,}680$ plantas y \$2.65~millones, 593.8 plantas/ha. Con las 128.8 plantas/ha existentes, la densidad al plantar es de 722.6 plantas/ha, 9.8\,\% sobre la meta: es la sobreplantación que absorbe la mortalidad del primer año con la confiabilidad exigida. El costo esperado de reposición es de apenas \$5/ha, porque con $\alpha=0.05$ el déficit esperado es casi nulo y el costo es casi todo de adquisición.

\paragraph{Composición} La composición óptima agota la banda \eqref{eq:banda} en la dirección que anticipaba el costo por planta sobreviviente $c_s/q_s$ (\cref{tab:compra-parametros}). Las cuatro especies de \textit{Opuntia}, las de menor $c_s/q_s$, suben entre 9.4\,\% y 9.9\,\% respecto al plan; \textit{P.~laevigata}, \textit{A.~scabra} y \textit{Y.~filifera}, las de mayor, bajan entre 9.2\,\% y 9.9\,\%; \textit{A.~lechuguilla} y \textit{A.~striata} bajan 8\,\%, y \textit{A.~salmiana} se mantiene. El margen de ajuste por especie, que el plan concibe como tolerancia de campo, resulta ser la principal palanca de costo del modelo.

\paragraph{Tamaño del polígono} Los dos polígonos más chicos son los más caros por hectárea (\$14{,}934 en el polígono 1, de 1.28~ha, y \$14{,}819 en el 4, de 1.38~ha), y la correlación entre superficie y $B_i^{*}$ es de $-0.63$. El margen necesario para cumplir con probabilidad $1-\alpha$ crece aproximadamente con la raíz de la meta, más lento que la meta misma, de modo que pesa más por hectárea cuando la meta es chica. La relación no es estricta porque $B_i^{*}$ también depende de cuántas plantas existentes tiene cada polígono.

\subsubsection{Distribución espacial}
\label{subsec:res-espacial}

\paragraph{Seis claves para 30 polígonos} Al repartir la compra sin presupuesto de cada polígono en 52 posiciones, los 30 polígonos se reducen a seis composiciones distintas y, por tanto, a seis claves de plantación (\cref{tab:res-mosaicos}). La primera la comparten 18 polígonos y la segunda, 5. En la práctica, la cuadrilla necesita solo seis tablas de claves para todo el sitio.

\begin{table}[t]
\centering
\caption{Claves de plantación. $n_p$: polígonos que comparten la clave; azar: competencia esperada de una asignación aleatoria con la misma composición, \cref{eq:competencia-azar}; memético: mejor competencia de 30 corridas; reducción: $1-C/\mathbb{E}[C]$.}
\label{tab:res-mosaicos}
\small
\begin{tabular}{@{}crrrr@{}}
\toprule
Clave & $n_p$ & Azar & Memético & Reducción (\%) \\
\midrule
1 & 18 & 93.34 & 78.18 & 16.2 \\
2 &  5 & 92.37 & 77.37 & 16.2 \\
3 &  1 & 95.09 & 81.29 & 14.5 \\
4 &  1 & 91.74 & 76.77 & 16.3 \\
5 &  4 & 91.73 & 76.80 & 16.3 \\
6 &  1 & 95.90 & 81.85 & 14.6 \\
\bottomrule
\end{tabular}
\end{table}

\paragraph{Reducción de la competencia} El memético reduce la competencia entre 14.5\,\% y 16.3\,\% respecto a una asignación aleatoria de igual composición. Las claves 3 y 6, de un solo polígono cada una, reducen menos porque tienen 6 posiciones de \textit{P.~laevigata} en lugar de 7, es decir, un separador menos, como se explica enseguida.

\paragraph{Estructura de las claves} La \cref{fig:clave1} muestra la clave 1. De los 156 pares de vecinas de cada mosaico, solo uno o dos son de la misma especie, y en todos los casos son de \textit{A.~salmiana}, que ocupa entre 27\,\% y 29\,\% de las posiciones. \textit{P.~laevigata}, con la mayor penalización intraespecífica (1.50), nunca queda junto a otra de su especie y actúa como separador: entre 45\,\% y 50\,\% de sus vecinas son \textit{A.~salmiana} y entre 29\,\% y 39\,\% \textit{O.~streptacantha}, las dos especies abundantes con penalización intraespecífica alta, con las que la suya es mínima (0.20 y 0.17). Es justo el arreglo de nodriza que motivó el término $F_{ab}$ de la matriz: cada mezquite queda rodeado de suculentas.

\paragraph{Comparación con la clave vigente} El orden de plantación que usa CONAFOR es una secuencia de 22 posiciones para cada 100~m lineales, con la composición del plan (siete \textit{A.~salmiana}, tres \textit{P.~laevigata} y entre una y dos de cada especie restante). Al ser una sola fila, la secuencia no dice cómo se acomodan las filas entre sí, y en el tresbolillo cada planta tiene dos vecinas en la fila anterior y dos en la siguiente. Se evaluó la secuencia en la retícula con la matriz $W$ para los 22 desfases posibles entre filas consecutivas (\cref{tab:clave-vigente}). Si cada fila repite la secuencia alineada con la anterior, cada planta queda junto a dos de su misma especie y la competencia es 13\,\% \emph{mayor} que la de una asignación al azar con la misma composición (coeficiente 1.13). Con el mejor desfase, de cinco posiciones, el coeficiente baja a 0.85, y el promedio de los 22 desfases es 0.94. La clave 1 del memético (coeficiente 0.84) mejora a la vigente en cualquier desfase, aunque frente al mejor de ellos la diferencia es pequeña; por par de vecinas, su competencia es 29\,\% menor que la de la secuencia alineada y 6\,\% menor que la del mejor desfase, aunque esta comparación mezcla dos composiciones distintas, la del plan y la de la compra. La mayor ganancia, por tanto, no viene solo de optimizar, sino de especificar el arreglo en dos dimensiones: una tabla de claves fija las vecinas de cada planta en todas las direcciones, mientras que una secuencia por fila deja la mitad de las vecindades al azar del desfase con que se empiece cada fila.

\begin{table}[t]
\centering
\caption{Clave vigente de 22 posiciones evaluada en el tresbolillo según el desfase entre filas consecutivas, frente a la clave 1 del memético. Coeficiente: competencia entre la esperada al azar con la misma composición; por par: competencia media por par de vecinas; mismas: vecinas de la misma especie por planta.}
\label{tab:clave-vigente}
\small
\begin{tabular}{@{}lrrr@{}}
\toprule
Arreglo & Coeficiente & Por par & Mismas \\
\midrule
Vigente, filas alineadas & 1.126 & 0.709 & 2.00 \\
Vigente, mejor desfase (5) & 0.850 & 0.535 & 0.09 \\
Vigente, media de los 22 desfases & 0.944 & 0.594 & 0.64 \\
Memético, clave 1 & 0.838 & 0.501 & 0.08 \\
\bottomrule
\end{tabular}
\end{table}

\begin{figure}
\centering
% Generado por scripts/figura_clave.py; no editar a mano.
% Mosaico 1: competencia 78.18, al azar 93.3365.
\begin{tikzpicture}[x=0.5cm,y=0.5cm,font=\sffamily\fontsize{4.5}{5}\selectfont]
  \filldraw[fill=green!45,draw=white,line width=0.4pt] (0.000,0.577) -- (-0.500,0.289) -- (-0.500,-0.289) -- (-0.000,-0.577) -- (0.500,-0.289) -- (0.500,0.289) -- cycle;
  \node at (0.000,0.000) {Asa};
  \filldraw[fill=brown!70,draw=white,line width=0.4pt] (1.000,0.577) -- (0.500,0.289) -- (0.500,-0.289) -- (1.000,-0.577) -- (1.500,-0.289) -- (1.500,0.289) -- cycle;
  \node at (1.000,0.000) {Pla};
  \filldraw[fill=orange!80!red,draw=white,line width=0.4pt] (2.000,0.577) -- (1.500,0.289) -- (1.500,-0.289) -- (2.000,-0.577) -- (2.500,-0.289) -- (2.500,0.289) -- cycle;
  \node at (2.000,0.000) {Ost};
  \filldraw[fill=green!45,draw=white,line width=0.4pt] (3.000,0.577) -- (2.500,0.289) -- (2.500,-0.289) -- (3.000,-0.577) -- (3.500,-0.289) -- (3.500,0.289) -- cycle;
  \node at (3.000,0.000) {Asa};
  \filldraw[fill=orange!20,draw=white,line width=0.4pt] (4.000,0.577) -- (3.500,0.289) -- (3.500,-0.289) -- (4.000,-0.577) -- (4.500,-0.289) -- (4.500,0.289) -- cycle;
  \node at (4.000,0.000) {Oca};
  \filldraw[fill=green!45,draw=white,line width=0.4pt] (5.000,0.577) -- (4.500,0.289) -- (4.500,-0.289) -- (5.000,-0.577) -- (5.500,-0.289) -- (5.500,0.289) -- cycle;
  \node at (5.000,0.000) {Asa};
  \filldraw[fill=brown!70,draw=white,line width=0.4pt] (6.000,0.577) -- (5.500,0.289) -- (5.500,-0.289) -- (6.000,-0.577) -- (6.500,-0.289) -- (6.500,0.289) -- cycle;
  \node at (6.000,0.000) {Pla};
  \filldraw[fill=orange!80!red,draw=white,line width=0.4pt] (7.000,0.577) -- (6.500,0.289) -- (6.500,-0.289) -- (7.000,-0.577) -- (7.500,-0.289) -- (7.500,0.289) -- cycle;
  \node at (7.000,0.000) {Ost};
  \filldraw[fill=green!45,draw=white,line width=0.4pt] (8.000,0.577) -- (7.500,0.289) -- (7.500,-0.289) -- (8.000,-0.577) -- (8.500,-0.289) -- (8.500,0.289) -- cycle;
  \node at (8.000,0.000) {Asa};
  \filldraw[fill=orange!60,draw=white,line width=0.4pt] (9.000,0.577) -- (8.500,0.289) -- (8.500,-0.289) -- (9.000,-0.577) -- (9.500,-0.289) -- (9.500,0.289) -- cycle;
  \node at (9.000,0.000) {Oro};
  \filldraw[fill=green!45,draw=white,line width=0.4pt] (10.000,0.577) -- (9.500,0.289) -- (9.500,-0.289) -- (10.000,-0.577) -- (10.500,-0.289) -- (10.500,0.289) -- cycle;
  \node at (10.000,0.000) {Asa};
  \filldraw[fill=green!10,draw=white,line width=0.4pt] (11.000,0.577) -- (10.500,0.289) -- (10.500,-0.289) -- (11.000,-0.577) -- (11.500,-0.289) -- (11.500,0.289) -- cycle;
  \node at (11.000,0.000) {Ast};
  \filldraw[fill=orange!60,draw=white,line width=0.4pt] (12.000,0.577) -- (11.500,0.289) -- (11.500,-0.289) -- (12.000,-0.577) -- (12.500,-0.289) -- (12.500,0.289) -- cycle;
  \node at (12.000,0.000) {Oro};
  \filldraw[fill=orange!40,draw=white,line width=0.4pt] (0.500,-0.289) -- (-0.000,-0.577) -- (0.000,-1.155) -- (0.500,-1.443) -- (1.000,-1.155) -- (1.000,-0.577) -- cycle;
  \node at (0.500,-0.866) {Oen};
  \filldraw[fill=green!45,draw=white,line width=0.4pt] (1.500,-0.289) -- (1.000,-0.577) -- (1.000,-1.155) -- (1.500,-1.443) -- (2.000,-1.155) -- (2.000,-0.577) -- cycle;
  \node at (1.500,-0.866) {Asa};
  \filldraw[fill=green!65!black!60,draw=white,line width=0.4pt] (2.500,-0.289) -- (2.000,-0.577) -- (2.000,-1.155) -- (2.500,-1.443) -- (3.000,-1.155) -- (3.000,-0.577) -- cycle;
  \node at (2.500,-0.866) {Asc};
  \filldraw[fill=green!45,draw=white,line width=0.4pt] (3.500,-0.289) -- (3.000,-0.577) -- (3.000,-1.155) -- (3.500,-1.443) -- (4.000,-1.155) -- (4.000,-0.577) -- cycle;
  \node at (3.500,-0.866) {Asa};
  \filldraw[fill=brown!70,draw=white,line width=0.4pt] (4.500,-0.289) -- (4.000,-0.577) -- (4.000,-1.155) -- (4.500,-1.443) -- (5.000,-1.155) -- (5.000,-0.577) -- cycle;
  \node at (4.500,-0.866) {Pla};
  \filldraw[fill=orange!80!red,draw=white,line width=0.4pt] (5.500,-0.289) -- (5.000,-0.577) -- (5.000,-1.155) -- (5.500,-1.443) -- (6.000,-1.155) -- (6.000,-0.577) -- cycle;
  \node at (5.500,-0.866) {Ost};
  \filldraw[fill=green!45,draw=white,line width=0.4pt] (6.500,-0.289) -- (6.000,-0.577) -- (6.000,-1.155) -- (6.500,-1.443) -- (7.000,-1.155) -- (7.000,-0.577) -- cycle;
  \node at (6.500,-0.866) {Asa};
  \filldraw[fill=violet!40,draw=white,line width=0.4pt] (7.500,-0.289) -- (7.000,-0.577) -- (7.000,-1.155) -- (7.500,-1.443) -- (8.000,-1.155) -- (8.000,-0.577) -- cycle;
  \node at (7.500,-0.866) {Yfi};
  \filldraw[fill=green!65!black!60,draw=white,line width=0.4pt] (8.500,-0.289) -- (8.000,-0.577) -- (8.000,-1.155) -- (8.500,-1.443) -- (9.000,-1.155) -- (9.000,-0.577) -- cycle;
  \node at (8.500,-0.866) {Asc};
  \filldraw[fill=green!20,draw=white,line width=0.4pt] (9.500,-0.289) -- (9.000,-0.577) -- (9.000,-1.155) -- (9.500,-1.443) -- (10.000,-1.155) -- (10.000,-0.577) -- cycle;
  \node at (9.500,-0.866) {Ale};
  \filldraw[fill=orange!60,draw=white,line width=0.4pt] (10.500,-0.289) -- (10.000,-0.577) -- (10.000,-1.155) -- (10.500,-1.443) -- (11.000,-1.155) -- (11.000,-0.577) -- cycle;
  \node at (10.500,-0.866) {Oro};
  \filldraw[fill=green!45,draw=white,line width=0.4pt] (11.500,-0.289) -- (11.000,-0.577) -- (11.000,-1.155) -- (11.500,-1.443) -- (12.000,-1.155) -- (12.000,-0.577) -- cycle;
  \node at (11.500,-0.866) {Asa};
  \filldraw[fill=green!10,draw=white,line width=0.4pt] (12.500,-0.289) -- (12.000,-0.577) -- (12.000,-1.155) -- (12.500,-1.443) -- (13.000,-1.155) -- (13.000,-0.577) -- cycle;
  \node at (12.500,-0.866) {Ast};
  \filldraw[fill=orange!60,draw=white,line width=0.4pt] (1.000,-1.155) -- (0.500,-1.443) -- (0.500,-2.021) -- (1.000,-2.309) -- (1.500,-2.021) -- (1.500,-1.443) -- cycle;
  \node at (1.000,-1.732) {Oro};
  \filldraw[fill=orange!40,draw=white,line width=0.4pt] (2.000,-1.155) -- (1.500,-1.443) -- (1.500,-2.021) -- (2.000,-2.309) -- (2.500,-2.021) -- (2.500,-1.443) -- cycle;
  \node at (2.000,-1.732) {Oen};
  \filldraw[fill=brown!70,draw=white,line width=0.4pt] (3.000,-1.155) -- (2.500,-1.443) -- (2.500,-2.021) -- (3.000,-2.309) -- (3.500,-2.021) -- (3.500,-1.443) -- cycle;
  \node at (3.000,-1.732) {Pla};
  \filldraw[fill=orange!80!red,draw=white,line width=0.4pt] (4.000,-1.155) -- (3.500,-1.443) -- (3.500,-2.021) -- (4.000,-2.309) -- (4.500,-2.021) -- (4.500,-1.443) -- cycle;
  \node at (4.000,-1.732) {Ost};
  \filldraw[fill=green!45,draw=white,line width=0.4pt] (5.000,-1.155) -- (4.500,-1.443) -- (4.500,-2.021) -- (5.000,-2.309) -- (5.500,-2.021) -- (5.500,-1.443) -- cycle;
  \node at (5.000,-1.732) {Asa};
  \filldraw[fill=brown!70,draw=white,line width=0.4pt] (6.000,-1.155) -- (5.500,-1.443) -- (5.500,-2.021) -- (6.000,-2.309) -- (6.500,-2.021) -- (6.500,-1.443) -- cycle;
  \node at (6.000,-1.732) {Pla};
  \filldraw[fill=orange!40,draw=white,line width=0.4pt] (7.000,-1.155) -- (6.500,-1.443) -- (6.500,-2.021) -- (7.000,-2.309) -- (7.500,-2.021) -- (7.500,-1.443) -- cycle;
  \node at (7.000,-1.732) {Oen};
  \filldraw[fill=orange!60,draw=white,line width=0.4pt] (8.000,-1.155) -- (7.500,-1.443) -- (7.500,-2.021) -- (8.000,-2.309) -- (8.500,-2.021) -- (8.500,-1.443) -- cycle;
  \node at (8.000,-1.732) {Oro};
  \filldraw[fill=orange!20,draw=white,line width=0.4pt] (9.000,-1.155) -- (8.500,-1.443) -- (8.500,-2.021) -- (9.000,-2.309) -- (9.500,-2.021) -- (9.500,-1.443) -- cycle;
  \node at (9.000,-1.732) {Oca};
  \filldraw[fill=green!45,draw=white,line width=0.4pt] (10.000,-1.155) -- (9.500,-1.443) -- (9.500,-2.021) -- (10.000,-2.309) -- (10.500,-2.021) -- (10.500,-1.443) -- cycle;
  \node at (10.000,-1.732) {Asa};
  \filldraw[fill=green!10,draw=white,line width=0.4pt] (11.000,-1.155) -- (10.500,-1.443) -- (10.500,-2.021) -- (11.000,-2.309) -- (11.500,-2.021) -- (11.500,-1.443) -- cycle;
  \node at (11.000,-1.732) {Ast};
  \filldraw[fill=green!65!black!60,draw=white,line width=0.4pt] (12.000,-1.155) -- (11.500,-1.443) -- (11.500,-2.021) -- (12.000,-2.309) -- (12.500,-2.021) -- (12.500,-1.443) -- cycle;
  \node at (12.000,-1.732) {Asc};
  \filldraw[fill=green!20,draw=white,line width=0.4pt] (13.000,-1.155) -- (12.500,-1.443) -- (12.500,-2.021) -- (13.000,-2.309) -- (13.500,-2.021) -- (13.500,-1.443) -- cycle;
  \node at (13.000,-1.732) {Ale};
  \filldraw[fill=violet!40,draw=white,line width=0.4pt] (1.500,-2.021) -- (1.000,-2.309) -- (1.000,-2.887) -- (1.500,-3.175) -- (2.000,-2.887) -- (2.000,-2.309) -- cycle;
  \node at (1.500,-2.598) {Yfi};
  \filldraw[fill=orange!80!red,draw=white,line width=0.4pt] (2.500,-2.021) -- (2.000,-2.309) -- (2.000,-2.887) -- (2.500,-3.175) -- (3.000,-2.887) -- (3.000,-2.309) -- cycle;
  \node at (2.500,-2.598) {Ost};
  \filldraw[fill=green!45,draw=white,line width=0.4pt] (3.500,-2.021) -- (3.000,-2.309) -- (3.000,-2.887) -- (3.500,-3.175) -- (4.000,-2.887) -- (4.000,-2.309) -- cycle;
  \node at (3.500,-2.598) {Asa};
  \filldraw[fill=brown!70,draw=white,line width=0.4pt] (4.500,-2.021) -- (4.000,-2.309) -- (4.000,-2.887) -- (4.500,-3.175) -- (5.000,-2.887) -- (5.000,-2.309) -- cycle;
  \node at (4.500,-2.598) {Pla};
  \filldraw[fill=orange!80!red,draw=white,line width=0.4pt] (5.500,-2.021) -- (5.000,-2.309) -- (5.000,-2.887) -- (5.500,-3.175) -- (6.000,-2.887) -- (6.000,-2.309) -- cycle;
  \node at (5.500,-2.598) {Ost};
  \filldraw[fill=green!45,draw=white,line width=0.4pt] (6.500,-2.021) -- (6.000,-2.309) -- (6.000,-2.887) -- (6.500,-3.175) -- (7.000,-2.887) -- (7.000,-2.309) -- cycle;
  \node at (6.500,-2.598) {Asa};
  \filldraw[fill=orange!20,draw=white,line width=0.4pt] (7.500,-2.021) -- (7.000,-2.309) -- (7.000,-2.887) -- (7.500,-3.175) -- (8.000,-2.887) -- (8.000,-2.309) -- cycle;
  \node at (7.500,-2.598) {Oca};
  \filldraw[fill=green!45,draw=white,line width=0.4pt] (8.500,-2.021) -- (8.000,-2.309) -- (8.000,-2.887) -- (8.500,-3.175) -- (9.000,-2.887) -- (9.000,-2.309) -- cycle;
  \node at (8.500,-2.598) {Asa};
  \filldraw[fill=brown!70,draw=white,line width=0.4pt] (9.500,-2.021) -- (9.000,-2.309) -- (9.000,-2.887) -- (9.500,-3.175) -- (10.000,-2.887) -- (10.000,-2.309) -- cycle;
  \node at (9.500,-2.598) {Pla};
  \filldraw[fill=orange!20,draw=white,line width=0.4pt] (10.500,-2.021) -- (10.000,-2.309) -- (10.000,-2.887) -- (10.500,-3.175) -- (11.000,-2.887) -- (11.000,-2.309) -- cycle;
  \node at (10.500,-2.598) {Oca};
  \filldraw[fill=green!20,draw=white,line width=0.4pt] (11.500,-2.021) -- (11.000,-2.309) -- (11.000,-2.887) -- (11.500,-3.175) -- (12.000,-2.887) -- (12.000,-2.309) -- cycle;
  \node at (11.500,-2.598) {Ale};
  \filldraw[fill=orange!60,draw=white,line width=0.4pt] (12.500,-2.021) -- (12.000,-2.309) -- (12.000,-2.887) -- (12.500,-3.175) -- (13.000,-2.887) -- (13.000,-2.309) -- cycle;
  \node at (12.500,-2.598) {Oro};
  \filldraw[fill=green!45,draw=white,line width=0.4pt] (13.500,-2.021) -- (13.000,-2.309) -- (13.000,-2.887) -- (13.500,-3.175) -- (14.000,-2.887) -- (14.000,-2.309) -- cycle;
  \node at (13.500,-2.598) {Asa};
\end{tikzpicture}

\caption{Clave de plantación 1, compartida por 18 polígonos: mosaico de $4\times13$ posiciones al tresbolillo que se repite sobre el polígono; competencia 78.18 frente a 93.34 al azar. Verdes: \textit{Agave}; naranjas: \textit{Opuntia}; café: \textit{P.~laevigata}; morado: \textit{Y.~filifera}. Abreviaturas como en la \cref{tab:matriz-competencia}.}
\label{fig:clave1}
\end{figure}

\subsubsection{Sensibilidad a la matriz de competencia}
\label{subsec:res-sensibilidad}

La matriz $W$ depende de rasgos aproximados y de dos parámetros fijados por supuesto. Para medir cuánto dependen las claves de esos valores, se despejó el producto $T_{ab}N_{ab}$ de cada par a partir de la matriz base y de \eqref{eq:matriz-W}, con $\gamma=1.5$, $\phi=0.3$ y el efecto nodriza aplicado a los pares de \textit{P.~laevigata} con agaves y nopales\footnote{Con esa convención la reconstrucción devuelve exactamente la matriz base. Aplicarlo también a \textit{Y.~filifera} implicaría un traslape de raíces con el mezquite de 0.75, mayor que el de cualquier suculenta, lo que no corresponde a una roseta de raíz somera.}. Con ello se construyeron 45 matrices alternativas: 25 en la rejilla $\gamma\in\{1,1.25,1.5,1.75,2\}$, $\phi\in\{0,0.15,0.3,0.45,0.6\}$, y 20 en las que el producto $T_{ab}N_{ab}$ de cada par se multiplica por un factor uniforme e independiente en $[0.7,1.3]$, una perturbación más severa que variar el diámetro de copa de cada especie, porque afecta a cada par por separado. Con cada matriz se reoptimizaron los seis mosaicos (10 corridas del memético) y se comparó la competencia de la clave base, evaluada con la matriz alternativa, con la de la clave reoptimizada.

Las claves son robustas a $\gamma$ y $\phi$: en las 150 combinaciones de matriz y mosaico de la rejilla, la clave base queda a lo más 1.6\,\% por encima de la reoptimizada (mediana 0.15\,\%), y con $\gamma\geq1.25$ y $\phi\geq0.3$ es igual de buena. Solo con $\gamma=1$, es decir, sin penalización intraespecífica adicional, la diferencia supera el 1\,\%. Los rasgos pesan más. Con las matrices perturbadas, cuya correlación de Spearman con la base va de 0.91 a 0.95, la clave base queda en promedio 6.6\,\% por encima de la reoptimizada (máximo 11.6\,\%), aunque sigue reduciendo la competencia entre 9\,\% y 20\,\% respecto al azar. En la práctica, $\gamma$ y $\phi$ no necesitan precisarse más, pero vale la pena afinar los diámetros de copa y los perfiles de raíz antes de imprimir las claves.

\subsection{Discusión}
\label{subsec:discusion}

\paragraph{Lo que significa para CONAFOR} Con el presupuesto actual no es posible cumplir la meta en ningún polígono con una confiabilidad razonable; se necesitan alrededor de \$14{,}500/ha, o bien reducir la meta o la superficie. El modelo traduce esa decisión a números: \cref{tab:frontera} indica cuánto cuesta cada nivel de riesgo, y la diferencia entre niveles es pequeña frente al faltante. Lo que sí está en manos del planeador es la composición: agotar el margen de $\pm$10\,\% hacia los nopales es lo que hace la compra más barata. Si ese margen se redujera por razones ecológicas, el presupuesto necesario subiría.

\paragraph{Limitaciones} Los resultados dependen de supuestos que conviene revisar con el socio formador. La supervivencia se estimó con un solo año de monitoreo y se supone independiente entre plantas; un año seco afectaría a todas a la vez. Las plantas existentes se toman como un conteo exacto que no muere. La matriz de competencia se construyó con rasgos aproximados; las claves son robustas a $\gamma$ y $\phi$, pero no a errores grandes en los rasgos (\cref{subsec:res-sensibilidad}), y la matriz no se ha contrastado con observaciones de campo. Por último, el agua, la disponibilidad en vivero y el terreno de cada polígono quedaron fuera del modelo.


\section{Conclusiones y recomendaciones}
\label{sec:conclusiones}

Este trabajo propuso un modelo de dos niveles para la reforestación de 30 polígonos del matorral desértico de San Luis Potosí. El primero decide cuántas plantas de cada especie comprar en cada polígono para alcanzar la meta de 658 plantas por hectárea con una probabilidad prefijada, al menor costo y dentro del presupuesto. El segundo acomoda esas plantas en el tresbolillo para minimizar la competencia entre vecinas. Ambos se resolvieron con un método aproximado validado contra uno exacto, y se implementaron en una herramienta reproducible que CONAFOR puede usar con sus propios datos.

\paragraph{Compra} La restricción de probabilidad se reduce al cuantil de una binomial negativa, que la simulación Monte Carlo estima con un error del orden de \$1/ha frente al cálculo exacto. El resultado principal es que el presupuesto de \$12{,}500/ha no alcanza en ningún polígono: cumplir la meta con 95\,\% de confianza cuesta en promedio \$14{,}496/ha, 16\,\% más, es decir, \$365~mil adicionales para los 183.01~ha. La infactibilidad es estructural: el solo costo de cubrir el faltante, corregido por mortalidad, ya excede el presupuesto. Aceptar más riesgo apenas ayuda, pues pasar de 1\,\% a 50\,\% de probabilidad de incumplimiento ahorra 3.7\,\%. El margen de ajuste por especie, en cambio, es la principal palanca de costo: la compra óptima lo agota hacia los nopales, las especies con menor costo por planta sobreviviente. La compra resultante sobreplanta 9.8\,\% sobre la meta para absorber la mortalidad del primer año.

\paragraph{Distribución} Al redondear la compra a un mosaico de 52 posiciones, los 30 polígonos se reducen a seis claves de plantación. El algoritmo memético reduce la competencia entre 14.5\,\% y 16.3\,\% respecto al azar, con una variación casi nula entre corridas. El modelo entero no encontró una clave mejor en ningún caso y acota su distancia al óptimo en 4.5\,\% como máximo, una cota que refleja sobre todo la debilidad de la relajación lineal. Las claves eliminan casi toda vecindad entre plantas de la misma especie y usan al mezquite como separador y nodriza de las suculentas. Frente al orden vigente, la ganancia más grande no viene solo de optimizar, sino de dar el arreglo en dos dimensiones: la secuencia actual de 22 posiciones, repetida fila tras fila sin desfase, coloca a cada planta junto a dos de su especie y compite más que una plantación al azar. Las claves son robustas a los dos parámetros supuestos de la matriz de competencia, pero un error de $\pm$30\,\% en los rasgos de las especies las deja en promedio 6.6\,\% por encima de las que se obtendrían con los rasgos correctos.

\paragraph{Simulación de la siembra} En 1{,}000 siembras simuladas de una hectárea, el valor esperado de las plantas existentes es de 128.6 por hectárea, y para cumplir la meta de cada especie con 95\,\% de probabilidad hay que suministrar 651 plantas, con las que la hectárea queda en promedio con 722.4 plantas vivas al año. Cumplir cada especie con 95\,\% de probabilidad no equivale a cumplirlas todas: las diez a la vez se cumplen en 68.4\,\% de las réplicas. Con la clave 1, el valor esperado de la competencia al año es de 802.7 por hectárea, 18.0\,\% menos que una siembra al azar y 35.8\,\% menos que la secuencia vigente con filas alineadas. La matriz de transición de la clave tiene la diagonal nula, igual que la vigente; la diferencia entre ambas está en las vecinas de las filas contiguas, que solo una clave en dos dimensiones controla.

\paragraph{Recomendaciones} Para CONAFOR se desprenden cinco recomendaciones: (i) ajustar el presupuesto a alrededor de \$14{,}500/ha, o reducir la meta o la superficie a lo que el presupuesto actual cubre; (ii) aprovechar el margen de ajuste por especie hacia los nopales; (iii) entregar a las cuadrillas tablas de claves en dos dimensiones en lugar de una secuencia por fila; (iv) si se requiere que todas las especies cumplan a la vez con 95\,\% de probabilidad, repartir el riesgo entre ellas con $\alpha=0.005$ por especie, que por la desigualdad de Boole garantiza al menos 95\,\% de cumplimiento conjunto y eleva el presupuesto mínimo de \$14{,}496 a \$14{,}728/ha, apenas 1.6\,\%; y (v) continuar el monitoreo de supervivencia por especie y registrar la posición de las plantas, que son los datos que más limitan al modelo.

\paragraph{Trabajo futuro} El modelo puede extenderse en cuatro direcciones. La primera es un horizonte de varios años, con la cadena de Markov por posición de la \cref{subsec:markov}, ya que el monitoreo muestra mortalidad también en el segundo año. La segunda es la supervivencia correlacionada entre plantas por años secos y la incertidumbre en los conteos de existentes, ambas incorporables sorteándolas por escenario. La tercera es calibrar la matriz de competencia con datos de campo: con las coordenadas de cada planta, las funciones de correlación de pares estimarían la interacción entre especies que aquí se construyó con rasgos. La cuarta es acoplar los dos niveles, porque la competencia que reduce la distribución espacial debería elevar la supervivencia y, con ella, abaratar la compra.


\printbibliography

\end{document}
